\documentclass[11pt,a4paper]{article}
\usepackage[T1]{fontenc}
\usepackage[utf8]{inputenc}
\usepackage{amsmath,amssymb,amsthm}
\usepackage{booktabs}
\usepackage{longtable}
\usepackage{geometry}
\geometry{margin=2.3cm}
\usepackage{hyperref}
\hypersetup{colorlinks=true,linkcolor=blue,citecolor=blue}

\newcommand{\E}{\mathbb{E}}
\newcommand{\Var}{\operatorname{Var}}
\newcommand{\Cov}{\operatorname{Cov}}
\newcommand{\dd}{\mathrm{d}}
\newcommand{\Tstop}{T_{\text{stop}}}
\newcommand{\Mrv}{M}
\newcommand{\icx}{\preceq_{\mathrm{icx}}}
\newcommand{\stle}{\preceq_{\mathrm{st}}}

\theoremstyle{plain}
\newtheorem{proposition}{Proposition}
\newtheorem{theorem}{Theorem}
\newtheorem{lemma}{Lemma}
\newtheorem{corollary}{Corollary}
\theoremstyle{remark}
\newtheorem{remark}{Remark}

\title{Extended evidence accumulation with motor preparation:\\
formal predictions, simulations, and the results that do not follow\\
\large Working document}
\date{}

\begin{document}
\maketitle

\section{Scope and summary}

This is a working document, not a draft for the manuscript.

It documents a formal implementation of the extended evidence accumulation theory, derives
its predictions, calibrates it against the J~Neuroscience experiment, and reports which
electromyographic results follow from it and which do not. Two variants are treated in
parallel, differing only in when motor preparation begins.

\paragraph{What works.} With preparation anchored to stimulus offset, sixteen of the
nineteen directional effects come out with the correct sign and are stable across simulation
seeds. Premotor execution time is reproduced to within fourteen milliseconds in every
condition, premature responses to within one percentage point, and motor execution time is
invariant as observed.

\paragraph{What does not.} R15, the duration by coherence interaction on delay-period
partial-activation rate, comes out with the opposite sign, and R17 fails with it as a
corollary. Section~\ref{sec:R15} shows this is not a calibration failure but a structural
incompatibility: R15 and the near-invariance of accuracy with stimulus duration constrain
the internal threshold in opposite directions. Nine mechanisms were tested; none removes the
contradiction.

\paragraph{A second unresolved point.} R3 is reproduced here but does not appear in Achard,
Gajdos Preuss \& Servant (2025), Experiment~2, which uses a nearly identical design.
Section~\ref{sec:R3} locates the difference precisely and reports that the explanation I
first proposed is refuted by the accuracy data.

\paragraph{Convention on statistics.} Effects reported as significant are the ANOVA results
of the manuscript. Additional analyses with no ANOVA counterpart, in particular
subject-level analyses of R15, use planned contrasts and are labelled as such.

%======================================================================
\section{The empirical results}
\label{sec:results}

The labels R1 to R24 are defined once here. They are exhaustive for electromyography: the
fifteen significant effects on partial-activation rate in Table~2, the four entries for the
two latencies in Table~1, and five statements in the Results and Discussion text.
Electroencephalographic measures are excluded.

\subsection{Latencies}

\begin{description}
\item[R1] Premotor execution time, the latency from the response signal to the onset of the
muscle activation producing the response, is shorter after a 2-second stimulus.
\item[R2] Premotor execution time is shorter at higher motion coherence.
\item[R3] Duration by coherence interaction on premotor execution time: the coherence effect
is larger in the 0.3-second condition.
\item[R4] Motor execution time is unaffected by any factor.
\end{description}

\subsection{Channel selectivity}

\begin{description}
\item[R5] The rate of partial muscle activations is near zero in the channel of the
non-selected response.
\end{description}

\subsection{Stimulus period, 0 to 2 s after stimulus onset}

\begin{description}
\item[R6] Main effect of duration. \item[R7] Main effect of coherence.
\item[R8] Main effect of time bin: the rate ramps upward.
\item[R9] Duration by time interaction: steeper ramp in the 0.3-second condition.
\item[R10] Coherence by time interaction: shallower ramp at lower coherence.
\item[R11] No duration by coherence interaction. This appears only as an absent cell in
Table~2 and is never discussed; listed for completeness.
\end{description}

\subsection{Delay period, 1 to 3 s after stimulus offset}

\begin{description}
\item[R12] Main effect of duration: higher rate after a 2-second stimulus.
\item[R13] Main effect of coherence. \item[R14] Main effect of time bin: continued ramping.
\item[R15] Duration by coherence interaction: the coherence effect is more pronounced after
a 2-second stimulus.
\item[R16] Duration by time interaction: the gap between duration conditions grows.
\item[R17] No coherence by time interaction, used as a diagnostic that evidence accumulation
has ceased.
\end{description}

\subsection{Response-signal period and response period}

\begin{description}
\item[R18] Main effect of time bin: the rate decreases. \item[R19] Main effect of duration.
\item[R20] Duration by time interaction. \item[R21] No coherence effects.
\item[R22] Main effect of time bin in the response period.
\item[R23] Duration by time interaction in the response period.
\item[R24] The rate appears to level off shortly before the response signal in the
2-second, high-coherence condition.
\end{description}


% =====================================================================
%  REPLACEMENT FOR SECTION 4 ("Theory") OF emg_results_classification.tex
%  Plus a new Section 4bis mapping theory to R1-R24.
%  Uses the same preamble (amsmath, amsthm, booktabs, \E \Var \Cov \dd \Tstop).
%  Additional macros required:
%     \newcommand{\Mrv}{M}
%     \newtheorem{theorem}{Theorem}
%     \newcommand{\icx}{\preceq_{\mathrm{icx}}}
%     \newcommand{\stle}{\preceq_{\mathrm{st}}}
% =====================================================================

\section{Theory}
\label{sec:theory}

The organising principle of this section is different from the previous version.
Rather than developing the stochastic machinery first and applying it afterwards,
every subsection below terminates in a statement about one or more of R1--R24, and
every such statement carries an explicit status:

\begin{description}
\item[\textsc{Structural}] true by construction of the model; reproducing it is not
evidence for the model.
\item[\textsc{Theorem}] follows from the model with no further assumption.
\item[\textsc{Conditional}] follows from the model given an explicitly stated property of
the accumulation stage, labelled (A1)--(A5), each of which is checkable by direct
simulation in one line.
\item[\textsc{Decomposition}] the sign is not determined; the effect is written as a sum of
terms of known sign, which yields a new testable prediction.
\item[\textsc{Not predicted}] outside the reach of the present formalism.
\end{description}

Section~\ref{sec:errata} lists the corrections relative to the previous version of this
section. Section~\ref{sec:numbers} collects the recomputed constants; several of them
differ materially from those previously quoted.

Throughout, the calibrated values are $\lambda=21.47$, $\kappa=10.778$,
$\sigma_m=1.8216$, $\alpha=0.1344$, $g=1.5897$, $r=0.6751$, $\theta=0.8931$, $w=6.16$,
$\sigma_q=1.213$, $T_{\text{enc}}=0.22$, all in the units fixed by $\sigma=1$.

%----------------------------------------------------------------------
\subsection{Standing assumptions}
\label{sec:assump}

\paragraph{(A0) Model equations.} On $[T_{\text{enc}},\Tstop]$,
\begin{equation}
\dd x=v_c\dd t+\sigma\dd W_x,\quad x(T_{\text{enc}})=0,
\qquad \dot{\tilde y}=\lambda(x-\tilde y),\quad \tilde y(T_{\text{enc}})=0,
\label{eq:acc}
\end{equation}
with $v_c=A(c/0.5)^{0.6}$ and
$\Tstop=\min\{\inf\{t:|\tilde y|\ge\theta\},\,T_{\text{enc}}+D\}$. Write
$T_\theta=\inf\{t:|\tilde y|\ge\theta\}$, $s=\operatorname{sign}\tilde y(\Tstop)$,
$m_0=\tilde y(\Tstop)$ and
\begin{equation}
\Mrv:=|m_0|\in[0,\theta].
\end{equation}
For $t>\Tstop$ the filter input is $m_0+\varepsilon$ with $\varepsilon(\Tstop)=0$ and
$\dd\varepsilon=-\kappa\varepsilon\dd t+\sigma_m\dd W_\varepsilon$. Preparation is added
outside the filter, anchored at $t_a=T_{\text{enc}}+D$ and frozen at
$t_R+T_{\text{enc}}$. From $t_R+T_{\text{enc}}$, $\dd q=w\dd t+\sigma_q\dd W_q$ adds
$s\,q^{+}$. The button is pressed when $z$ reaches $g+r$, plus $T_r$.

\paragraph{(A1) Independence of the maintenance noise.}
$W_\varepsilon\perp\!\!\!\perp W_x$, and $W_q\perp\!\!\!\perp(W_x,W_\varepsilon)$.

This assumption is doing real work and was previously left implicit. The old text
justified $sU\stackrel{d}{=}U$ ``because $\varepsilon$ is a symmetric process''. That is
not sufficient: $s$ is $\mathcal F^{x}_{\Tstop}$-measurable, so symmetry of the law of
$U$ alone does not licence dropping $s$ --- one needs $U$ to be independent of $s$. Under
(A1) it is, and the following holds exactly.

\begin{lemma}[Conditional representation]
\label{lem:cond}
Under (A0)--(A1), conditionally on $(\Tstop,s,m_0)$,
\begin{equation}
z(t)=\mu(t)+V(t-\Tstop),\qquad
\mu(t)=\Mrv+\alpha\bigl(\min(t,t_R+T_{\text{enc}})-t_a\bigr)^{+},
\label{eq:decomp}
\end{equation}
where $V$ is a centred Gaussian process, independent of $(\Tstop,s,m_0)$, with the law of
$U$ solving $\dot U=\lambda(\varepsilon-U)$, $U(0)=\varepsilon(0)=0$.
\end{lemma}

\begin{proof}
$U$ is a measurable functional of $\{W_\varepsilon(\Tstop+u)-W_\varepsilon(\Tstop):u\ge0\}$
alone, because $\varepsilon$ and $U$ are both reset to $0$ at $\Tstop$. By (A1) that
increment process is independent of $\mathcal F^{x}_{\Tstop}\supseteq\sigma(\Tstop,s,m_0)$
and is a Brownian motion. Hence $U\perp\!\!\!\perp(\Tstop,s,m_0)$, and since the law of
$U$ is symmetric, $sU\stackrel{d}{=}U$ conditionally on $s$.
\end{proof}

\begin{remark}
If instead the maintenance perturbation were the continuation of the accumulation noise,
Lemma~\ref{lem:cond} would fail: conditioning on $\{\tilde y \text{ hits }+\theta\}$ tilts
the recent increments of $W_x$ upward, so $sU$ would start with a positive bias of order
$\sigma\sqrt{\lambda/2}\cdot\lambda^{-1}$ decaying over $1/\lambda\approx47$~ms. The
present model avoids this by construction. It should be stated as a modelling choice
rather than derived from symmetry.
\end{remark}

\paragraph{(A2)--(A5): accumulation-stage properties.}
Everything below that concerns condition effects is a statement about the law of $\Mrv$
and of $T_\theta$. Four properties are used. Let
$F_{D,c}=\Pr(T_\theta\le T_{\text{enc}}+D)$.

\begin{description}
\item[(A2)] \emph{Duration monotonicity.} $F_{D_2,c}\ge F_{D_1,c}$ for $D_2>D_1$, and
$\Mrv_{D_1,c}\stle \Mrv_{D_2,c}$.
\item[(A3)] \emph{Coherence monotonicity.} $F_{D,c_H}\ge F_{D,c_L}$ and
$\Mrv_{D,c_L}\stle \Mrv_{D,c_H}$ for $c_H>c_L$.
\item[(A4)] \emph{Hitting-time monotonicity.} $T_\theta^{c_H}\stle T_\theta^{c_L}$.
\item[(A5)] \emph{Compression.} The coherence contrast in $\Mrv$ is smaller at long
duration: $\E\Mrv_{2,c_H}-\E\Mrv_{2,c_L}\le\E\Mrv_{0.3,c_H}-\E\Mrv_{0.3,c_L}$.
\end{description}

(A2) is immediate for $F$ because $T_\theta\le T_{\text{enc}}+D_1$ implies
$T_\theta\le T_{\text{enc}}+D_2$ pathwise. (A5) is a consequence of saturation at $\theta$
and is the property the manuscript invokes for R3. (A3)--(A4) are partly provable:

\begin{lemma}[Pathwise ordering of the stopped process]
\label{lem:coupling}
Couple two coherence levels on the same $W_x$ and write
$\hat y^{c}(t)=\tilde y^{c}(t\wedge T_\theta^{c})$. Then
$\hat y^{c_H}(t)\ge\hat y^{c_L}(t)$ for all $t$, almost surely. In particular
$m_0^{c_H}\ge m_0^{c_L}$ a.s.
\end{lemma}

\begin{proof}
Before either process stops,
$\tilde y^{c_H}(t)-\tilde y^{c_L}(t)=\Delta(t):=(v_{c_H}-v_{c_L})\bigl(t-(1-e^{-\lambda t})/\lambda\bigr)\ge0$,
since the two solutions of \eqref{eq:acc} driven by the same $W_x$ differ only by the
deterministic response to the drift difference. Suppose $H$ stops first. If it stops at
$-\theta$ then $\tilde y^{c_L}\le-\theta$ at that instant, so $L$ had already stopped, a
contradiction; hence $H$ stops at $+\theta$ and thereafter
$\hat y^{c_H}=\theta\ge\hat y^{c_L}$. Symmetrically, if $L$ stops first it stops at
$-\theta$ and thereafter $\hat y^{c_L}=-\theta\le\hat y^{c_H}$.
\end{proof}

\begin{lemma}[Likelihood-ratio ordering of the commitment time]
\label{lem:girsanov}
In the unfiltered limit $\lambda\to\infty$ (threshold applied to $x$), the density of
$T_\theta$ satisfies
$f_{v}(t)/f_{v'}(t)=\bigl[\cosh(v\theta/\sigma^2)/\cosh(v'\theta/\sigma^2)\bigr]
e^{-(v^2-v'^2)t/2\sigma^2}$, which is decreasing in $t$ for $v>v'$. Hence
$T_\theta^{c_H}\preceq_{\mathrm{lr}}T_\theta^{c_L}$ and (A4) holds.
For finite $\lambda$ the same Girsanov computation gives
$f_v(t)/f_0(t)=\E_0[\cosh(v x_{T_\theta}/\sigma^2)\mid T_\theta=t]\,e^{-v^2t/2\sigma^2}$,
and (A4) holds whenever the conditional expectation is non-increasing in $t$; the
departure is controlled by
$\operatorname{sd}(x-\tilde y)=\sigma/\sqrt{2\lambda}=0.153$ against $\theta=0.893$.
\end{lemma}

\begin{proof}
Girsanov gives $\dd P_v/\dd P_0|_{\mathcal F_{T_\theta}}
=\exp(vx_{T_\theta}/\sigma^2-v^2T_\theta/2\sigma^2)$. Under $P_0$ the system is invariant
under the sign flip $W\mapsto-W$, so conditionally on $T_\theta=t$ the law of
$x_{T_\theta}$ is symmetric and the exponential may be replaced by its even part.
When $\lambda=\infty$, $x_{T_\theta}=\pm\theta$ exactly and the conditional expectation is
the constant $\cosh(v\theta/\sigma^2)$.
\end{proof}

\begin{remark}[What Lemma~\ref{lem:coupling} does \emph{not} give]
The pathwise ordering is of the \emph{signed} $m_0$; the quantity that drives the muscle
is $\Mrv=|m_0|$, and $|\cdot|$ destroys the ordering on error trials. The obstruction is
confined to sub-threshold error trials, since $\Mrv=\theta$ whenever $T_\theta$ is reached,
irrespective of sign. (A3) is therefore exact whenever every trial reaches $\theta$, and
approximate with an error at most $\sup h\cdot\Pr(\Mrv<\theta,\ m_0<0)$ otherwise. It is
checkable in one line of simulation and should be checked rather than assumed.
\end{remark}

%----------------------------------------------------------------------
\subsection{The gating lemma: nothing crosses before commitment}
\label{sec:gate}

This is the single most useful structural fact in the model and it was absent from the
previous version.

\begin{lemma}[Gating]
\label{lem:gate}
Assume $\theta<g$ and offset anchoring. Then for every $t<\Tstop$, $z(t)<g$ almost surely.
Consequently every partial activation, in every analysis window, occurs on a trial that
has already committed, and
\begin{equation}
\nu_{D,c}(t)=\E\Bigl[\mathbf 1\{\Tstop\le t\}\;
\rho\bigl(t-\Tstop;\,g-\mu(t)\bigr)\Bigr],
\label{eq:master-stim}
\end{equation}
where $\rho(\tau;A)$ is the transient crossing rate of Proposition~\ref{prop:rice-gen}.
\end{lemma}

\begin{proof}
Before $\Tstop$, $|\tilde y|<\theta$ by definition of the first hitting time; under offset
anchoring $t_a=T_{\text{enc}}+D\ge\Tstop$, so the preparation term is zero on
$[0,\Tstop)$; hence $z=s\tilde y\le|\tilde y|<\theta<g$. Equation~\eqref{eq:master-stim}
is then Lemma~\ref{lem:cond} together with the tower property.
\end{proof}

$\theta<g$ is not a calibration accident: $\theta\ge g$ would make every committed trial
respond immediately. The lemma therefore holds for the whole admissible parameter region.

\begin{remark}[This corrects the previous Proposition~6]
The old argument for R9 read: ``the level rises at $\bar v$ for uncommitted trials and is
flat for committed ones''. The mean level of uncommitted trials is irrelevant, because
uncommitted trials contribute \emph{exactly zero} crossings. The stimulus-period rate is
governed entirely by the commitment-time distribution and by what happens after
commitment. Every stimulus-period statement below is rewritten accordingly.
\end{remark}

%----------------------------------------------------------------------
\subsection{The maintenance process: exact second moments}
\label{sec:cov}

\begin{lemma}[Kernel]
\label{lem:kernel}
$V(\tau)=\int_0^{\tau}K(\tau-p)\dd W(p)$ with
$K(v)=\dfrac{\lambda\sigma_m}{\lambda-\kappa}\bigl(e^{-\kappa v}-e^{-\lambda v}\bigr)$,
$K(0)=0$, $K'(0)=\lambda\sigma_m$.
\end{lemma}

The proof is unchanged. Two consequences that the previous version did not state
explicitly, and which are needed for the stimulus period, follow immediately from
$\dd/\dd\tau\int_0^\tau K^2=K(\tau)^2$:

\begin{proposition}[Exact transient moments]
\label{prop:transient}
For all $\tau\ge0$,
\begin{align}
\sigma_V^2(\tau)&=\int_0^\tau K^2
=\frac{\lambda^{2}\sigma_m^{2}}{(\lambda-\kappa)^{2}}
\left[\frac{1-e^{-2\kappa\tau}}{2\kappa}
-\frac{2(1-e^{-(\lambda+\kappa)\tau})}{\lambda+\kappa}
+\frac{1-e^{-2\lambda\tau}}{2\lambda}\right],\\
\sigma_{\dot V}^2(\tau)&=\int_0^\tau (K')^2
=\frac{\lambda^{2}\sigma_m^{2}}{(\lambda-\kappa)^{2}}
\left[\frac{\lambda(1-e^{-2\lambda\tau})}{2}
-\frac{2\lambda\kappa(1-e^{-(\lambda+\kappa)\tau})}{\lambda+\kappa}
+\frac{\kappa(1-e^{-2\kappa\tau})}{2}\right],\\
\Cov\bigl(V(\tau),\dot V(\tau)\bigr)&=\int_0^\tau KK'=\tfrac12 K(\tau)^2 .
\label{eq:covKK}
\end{align}
As $\tau\to\infty$ these converge to
$\sigma_Y^2=\lambda\sigma_m^2/(2\kappa(\lambda+\kappa))$,
$\sigma_{\dot Y}^2=\lambda^2\sigma_m^2/(2(\lambda+\kappa))$ and $0$, and
$C(\tau)=B[e^{-\kappa|\tau|}/\kappa-e^{-\lambda|\tau|}/\lambda]$ with
$B=\lambda^2\sigma_m^2/(2(\lambda-\kappa)(\lambda+\kappa))$;
$\beta:=\sigma_{\dot Y}/\sigma_Y=\sqrt{\lambda\kappa}$.
\end{proposition}

Identity \eqref{eq:covKK} is the clean form of the statement the old Remark~3 was
reaching for. The reduced Rice formula requires $\Cov(V,\dot V)=0$; this is
\emph{exactly} zero only in the limit, and the exact transient value is $K(\tau)^2/2$,
which decays at rate $2\kappa$. It is what must be carried in the stimulus period, where
$\tau$ can be a few tens of milliseconds.

\begin{proposition}[Rate of convergence, corrected]
\label{prop:conv}
$\dfrac{\sigma_Y^2-\sigma_V^2(\tau)}{\sigma_Y^2}
=\dfrac{\lambda(\lambda+\kappa)}{(\lambda-\kappa)^2}e^{-2\kappa\tau}\bigl(1+o(1)\bigr)
=6.056\,e^{-2\kappa\tau}\bigl(1+o(1)\bigr)$ at calibrated parameters.
\end{proposition}

The previous version quoted $e^{-2\kappa\tau}$, omitting the prefactor $6.06$. At
$\tau=1$~s the deficit is $2.6\times10^{-9}$ rather than $4\times10^{-10}$; nothing
empirical changes, but the stated bound was not correct.

\begin{proposition}[Path regularity]
\label{prop:C1}
$V\in C^1$ almost surely, and $V\notin C^2$ almost surely. Moreover
$C\in C^2(\mathbb R)\setminus C^3(\mathbb R)$.
\end{proposition}

\begin{proof}
Pathwise: $\varepsilon$ is an Ornstein--Uhlenbeck process, hence has continuous paths a.s.;
$U$ solves $\dot U=\lambda(\varepsilon-U)$, so $\dot U$ is a continuous function of two
continuous functions, hence continuous a.s.; therefore $U\in C^1$. If $U''$ existed then
$\dot\varepsilon=U''/\lambda+\dot U$ would exist, which fails a.s. for an
Ornstein--Uhlenbeck process. For $C$: on $\tau>0$,
$C'''(0^{+})=B(\lambda^2-\kappa^2)\neq0$, and $C'''$ is odd because $C$ is even, so
$C'''(0^{-})=-C'''(0^{+})$ and $C'''(0)$ does not exist.
\end{proof}

\begin{remark}[Why the previous proof of H1 was wrong]
The old text asserted that a stationary Gaussian process has a.s.\ continuously
differentiable paths ``whenever its covariance is twice differentiable at the origin with
$C''$ continuous''. That is false in general: continuity of the covariance of the
derivative process does not imply sample continuity of the derivative process, the
standard counterexample being $r(0)-r(\tau)\asymp1/\log(1/|\tau|)$. A sufficient condition
here would be $-C''(0)+C''(\tau)=O(|\tau|)$, which does hold, since
$-C''(0)+C''(\tau)=B[\lambda(1-e^{-\lambda|\tau|})-\kappa(1-e^{-\kappa|\tau|})]
=B(\lambda^2-\kappa^2)|\tau|+O(\tau^2)$. But the pathwise argument above is shorter and
needs no entropy condition, so the general theorem should simply be dropped.
The statement ``$C^2$ but not $C^4$'' should become ``$C\in C^2\setminus C^3$'', which is
strictly stronger and is proved in one line.
\end{remark}

%----------------------------------------------------------------------
\subsection{Kac--Rice in the form actually needed}
\label{sec:rice}

\begin{proposition}[General Gaussian crossing rate]
\label{prop:rice-gen}
Let $z(t)=\mu(t)+V(\tau)$ with $(V,\dot V)$ jointly Gaussian, $\Var V=s^2>0$,
$\Var\dot V=\dot s^2>0$, $\Cov(V,\dot V)=\gamma$, $\E\dot z=\bar\alpha$. Write
$A=g-\mu(t)$, $\eta=\bar\alpha+\gamma A/s^2$, $\omega^2=\dot s^2-\gamma^2/s^2$. Then the
upcrossing rate of $g$ is
\begin{equation}
\rho=\frac{1}{s}\varphi\!\Bigl(\frac{A}{s}\Bigr)
\Bigl[\eta\,\Phi(\eta/\omega)+\omega\,\varphi(\eta/\omega)\Bigr].
\label{eq:rice-gen}
\end{equation}
\end{proposition}

\begin{proof}
$\nu=p_{z}(g)\E[\dot z^{+}\mid z=g]$. Given $V=A$, $\dot V\sim\mathcal N(\gamma A/s^2,\omega^2)$,
so $\dot z\mid z=g\sim\mathcal N(\eta,\omega^2)$, and
$\E[\mathcal N(\eta,\omega^2)^{+}]=\eta\Phi(\eta/\omega)+\omega\varphi(\eta/\omega)$.
Hypotheses H1--H3 hold by Proposition~\ref{prop:C1} and Proposition~\ref{prop:transient}.
\end{proof}

\begin{corollary}[Delay-period reduction]
\label{cor:rice-delay}
In the delay window $\tau\ge1$~s, so $\gamma=0$ to within $10^{-9}$, $s=\sigma_Y$,
$\dot s=\sigma_{\dot Y}$, $\bar\alpha=\alpha$, and
\begin{equation}
\nu(t\mid \Mrv)=\Gamma_0\,h\bigl(\mu(t)\bigr),\qquad
h(u)=\frac{1}{\sigma_Y}\varphi\!\Bigl(\frac{g-u}{\sigma_Y}\Bigr),\qquad
\Gamma_0=\alpha\Phi(m)+\sigma_{\dot Y}\varphi(m),\ m=\frac{\alpha}{\sigma_{\dot Y}} .
\label{eq:nucond}
\end{equation}
$\Gamma_0=2.0107$ and $\nu_{\max}=\Gamma_0\varphi(0)/\sigma_Y=\Gamma\beta/2\pi=2.506$ per
second, where $\Gamma=\sqrt{2\pi}[m\Phi(m)+\varphi(m)]=1.0350$.
\end{corollary}

\begin{remark}[Corrected bound]
The old text bounded the rate by $\beta/2\pi=2.421$. With a trend the correct bound is
$\Gamma\beta/2\pi=2.506$. The distinction is $3.5$~per~cent and matters only for absolute
rates, but the R24 argument quotes the bound explicitly.
\end{remark}

Write $a(t)=(g-\mu(t))/\sigma_Y$ and $k=\alpha/\sigma_Y=0.4198$~s$^{-1}$. The derivatives
of $h$, used repeatedly below, are
\begin{equation}
h'=\frac{a}{\sigma_Y}h>0,\qquad
h''=\frac{a^2-1}{\sigma_Y^2}h,\qquad
h'''=\frac{a(a^2-3)}{\sigma_Y^{3}}h .
\label{eq:hderiv}
\end{equation}
So $h$ is increasing on $u<g$; convex where $a>1$; and $h'$ is convex where $a>\sqrt3$.
These three thresholds do all the work in Sections~\ref{sec:delaypred}
and~\ref{sec:diagnostics}, and the calibrated $a$ straddles $\sqrt3$
(Section~\ref{sec:numbers}).

%----------------------------------------------------------------------
\subsection{The three corrections}
\label{sec:corrections}

\subsubsection{Trend}
Unchanged: $\Gamma$ is a level-independent multiplicative factor of $1.0350$ that cancels
from every ratio of rates at the same time point and from every log-slope, and does not
cancel from absolute rates.

\subsubsection{Absorption}
The observed process is killed at $g+r$. The previous version asserted
$\nu^{\text{obs}}\le\nu$ on the ground that ``killing removes precisely the paths that have
been highest''. That reasoning is incomplete, and in one respect points the wrong way:

\begin{remark}[Two opposing effects]
Conditioning on survival lowers the density of $z(t)$ at $g$. But it also biases the slope
\emph{upward}, because $\Cov(V(u),\dot V(t))=C'(t-u)=B[-e^{-\kappa(t-u)}+e^{-\lambda(t-u)}]<0$
for $u<t$: the process is mean-reverting, so a path constrained to have been low is more
likely to be rising now. The net sign is not determined by inspection.
\end{remark}

\begin{proposition}[A bound that is actually proved]
\label{prop:absorb}
With $\tau_{\text{abs}}=\inf\{t:z(t)\ge g+r\}$ and $p(t)=\Pr(\tau_{\text{abs}}\le t)$,
\begin{equation}
\nu^{\text{obs}}(t)\le\frac{\nu(t)}{1-p(t)} .
\end{equation}
\end{proposition}

\begin{proof}
$\E[N_g(t,t+h)\mathbf 1\{\tau_{\text{abs}}>t\}]\le\E[N_g(t,t+h)]$; divide by $h$ and by
$\Pr(\tau_{\text{abs}}>t)=1-p(t)$ and let $h\downarrow0$.
\end{proof}

At the observed premature rate of $1.35$~per~cent the bound is $\nu^{\text{obs}}\le1.014\nu$.
This is weaker in form than the old claim but it is a theorem, and it is of the size the
simulation reports. \textbf{However}, the size of the correction quoted previously was
based on $a\in[2.2,3.1]$, which is wrong (Section~\ref{sec:numbers}): at
$a\in[1.4,1.8]$ the note's own absorption table gives a correction near $1$~per~cent, not
below $0.2$~per~cent.

\subsubsection{Discretisation}
\label{sec:discretisation}

\begin{proposition}[Euler bias is not multiplicative across conditions]
\label{prop:euler}
The Euler scheme with step $\Delta$ applied to $(\varepsilon,U)$ has stationary variance
\begin{equation}
\sigma_{Y,\Delta}^2=\sigma_Y^2\bigl(1+\delta\bigr)+O(\Delta^2),\qquad
\delta=\frac{\lambda\kappa}{\lambda+\kappa}\Delta=7.176\,\Delta .
\end{equation}
Consequently
\begin{equation}
\ln\nu_\Delta-\ln\nu=\frac{(a^2-1)\delta}{2}+\text{(terms independent of }a)+O(\Delta^2),
\end{equation}
so the relative bias differs between two conditions with distances $a_1,a_2$ by
$\exp[(a_1^2-a_2^2)\delta/2]$.
\end{proposition}

\begin{proof}
Let $A=1-\kappa\Delta$, $\rho=1-\lambda\Delta$, $b=\lambda\Delta$. Stationarity of the
discrete chain gives $\Var\varepsilon=\sigma_m^2\Delta/(1-A^2)$,
$\Cov(U,\varepsilon)=bA\Var\varepsilon/(1-\rho A)$ and
$\Var U(1-\rho^2)=b^2\Var\varepsilon+2\rho b\Cov(U,\varepsilon)$. Expanding to first order
in $\Delta$ gives $\Var\varepsilon=(\sigma_m^2/2\kappa)(1+\kappa\Delta/2)$ and
$\Var U=\frac{\lambda}{\lambda+\kappa}\Var\varepsilon\bigl[1+\tfrac{\kappa(\lambda-\kappa)}{2(\lambda+\kappa)}\Delta\bigr]$,
whose product is the stated $\delta$. The second display follows from
$\ln\nu=\text{const}-\ln\sigma_Y-a^2/2$ with $a^2\mapsto a^2/(1+\delta)$.
\end{proof}

At $\Delta=2$~ms, $\delta=0.0144$; across the observed spread $a\in[1.45,1.74]$ the
differential bias is $0.6$~per~cent. That is small in absolute terms but it is of the
order of thirty per cent of R16 and comparable with several seed standard deviations. The
statement in the previous version --- ``a multiplicative bias of this kind largely
cancels'' --- is therefore an empirical observation about the particular table, not a
property of the estimator.

\textbf{Recommendation.} $(\varepsilon,U)$ is a two-dimensional linear SDE with an exact
Gaussian transition kernel. Sampling it exactly removes $\delta$ entirely; detecting
crossings by quadratic interpolation between samples removes most of the residual
undercount. The convergence table in the previous version is not consistent with a single
$O(\Delta)$ source (the residual falls by less than a factor two when $\Delta$ falls by a
factor eight), which is a further reason to eliminate the leading term analytically rather
than extrapolate it numerically.

%----------------------------------------------------------------------
\subsection{The two master formulas}
\label{sec:master}

Everything about R6--R24 is a corollary of these two equations.

\begin{theorem}[Master formulas]
\label{thm:master}
Under (A0)--(A1) and Lemma~\ref{lem:gate}:
\begin{align}
\text{delay window }(\tau\ge1\text{ s}):&\quad
\boxed{\;\nu_{D,c}(t)=\Gamma_0\,\E\bigl[h\bigl(\Mrv_{D,c}+\alpha(t-t_a)\bigr)\bigr]\;}
\label{eq:master-delay}\\[4pt]
\text{stimulus window}:&\quad
\boxed{\;\nu_{D,c}(t)=\E\Bigl[\mathbf 1\{\Tstop\le t\}\,
\rho\bigl(t-\Tstop;\;g-\Mrv-\alpha(t-t_a)^{+}\bigr)\Bigr]\;}
\label{eq:master-stim2}
\end{align}
with $h$ as in \eqref{eq:nucond} and $\rho$ as in \eqref{eq:rice-gen}.
\end{theorem}

Equation~\eqref{eq:master-delay} says that \emph{every} delay-period effect is a
functional of the distribution of the maintained level $\Mrv$ alone. Equation
\eqref{eq:master-stim2} says that every stimulus-period effect is a functional of the
joint law of $(\Tstop,\Mrv)$ alone --- the commitment-time distribution convolved with a
fixed post-commitment response kernel.

For the stimulus period a further simplification applies. In the 2-second condition
$t_a=2.22$~s lies outside the analysis window, so preparation contributes nothing and
$\Mrv=\theta$ for every trial that contributes at all:
\begin{equation}
\nu_{2,c}(t)=\int_0^{t}f^{c}_{T_\theta}(u)\,\rho\bigl(t-u;\,g-\theta\bigr)\dd u .
\label{eq:stim2s}
\end{equation}
In the 0.3-second condition $t_a=0.52$~s lies \emph{inside} the window and every trial has
committed by then, at $\Mrv=\theta$ if $T_\theta\le0.52$ and at $\Mrv=|\tilde y(0.52)|<\theta$
otherwise.

%----------------------------------------------------------------------
\subsection{Predictions: delay period (R12--R17, R24)}
\label{sec:delaypred}

Throughout this subsection $a(t)>1$, which holds at calibrated parameters
(Section~\ref{sec:numbers}), so $h$ is increasing and convex on the relevant range.

\begin{theorem}[R14: the delay rate ramps --- unconditional]
\label{thm:R14}
If $\Mrv+\alpha(t-t_a)<g$ a.s.\ then $\partial_t\nu_{D,c}=\Gamma_0\alpha\E[h']>0$.
\emph{Status: \textsc{Theorem}.} At calibrated parameters
$\Mrv+3\alpha\le\theta+0.403=1.296<g=1.590$, so the hypothesis holds with margin.
\end{theorem}

\begin{theorem}[R12, R13: main effects of duration and coherence]
\label{thm:R1213}
If $\Mrv_{D_1,c}\stle\Mrv_{D_2,c}$ then $\nu_{D_1,c}(t)\le\nu_{D_2,c}(t)$ at every $t$
(R12); if $\Mrv_{D,c_L}\stle\Mrv_{D,c_H}$ then $\nu_{D,c_L}(t)\le\nu_{D,c_H}(t)$ (R13).
\emph{Status: \textsc{Conditional} on (A2), (A3).}
Only monotonicity of $h$ is used, hence only $\mu<g$ is required.
Moreover, decomposing on $\{T_\theta\le T_{\text{enc}}+D\}$,
\begin{equation}
\nu_{D,c}=\Gamma_0\Bigl[F_{D,c}h(\theta+\alpha\tau)
+(1-F_{D,c})\E\bigl[h(\Mrv+\alpha\tau)\mid \Mrv<\theta\bigr]\Bigr],
\label{eq:mixture}
\end{equation}
and because $h$ is convex the first term dominates: the condition effects are carried
mainly by $\Delta F$, not by the sub-threshold level.
\end{theorem}

\begin{theorem}[R16, and the proportionality law]
\label{thm:R16}
Under the same hypotheses,
$\partial_t\bigl[\nu_{D_2,c}-\nu_{D_1,c}\bigr]=\Gamma_0\alpha\bigl(\E[h'(\Mrv_{D_2}+\alpha\tau)]
-\E[h'(\Mrv_{D_1}+\alpha\tau)]\bigr)>0$ whenever $h'$ is increasing, i.e.\ whenever
$a>1$. Hence the duration gap grows: R16. \emph{Status: \textsc{Conditional} on (A2).}
More sharply, since $h'/h=a/\sigma_Y$,
\begin{equation}
\boxed{\;
\frac{\partial_t\bigl[\nu^{c_H}-\nu^{c_L}\bigr]}{\nu^{c_H}-\nu^{c_L}}
\in\Bigl[k\,a_{\min},\;k\,a_{\max}\Bigr]
=[0.59,\,0.76]\ \text{s}^{-1}\;}
\label{eq:proportionality}
\end{equation}
and identically for the duration contrast: \textbf{any} condition effect on the
delay-period rate that acts through the maintained level must grow at approximately
$0.6$--$0.75$ per second, i.e.\ roughly triple across the two-second analysis window.
\end{theorem}

\begin{corollary}[R17 is not a saturation null --- it contradicts R13]
\label{cor:R17}
R16 and R13 are both observed. Equation~\eqref{eq:proportionality} applies to the
coherence contrast exactly as it applies to the duration contrast. R17, the absence of a
coherence $\times$ time interaction, therefore asserts that the coherence contrast does
\emph{not} grow, while R13 asserts that it exists. These are jointly incompatible with any
account in which coherence acts on the delay-period rate through the maintained level ---
including but not restricted to the present model.
\emph{Status: \textsc{Theorem} (a constraint on the data, not on the model).}
\end{corollary}

This is a reclassification of R17. The previous version treated R17 as a corollary of R15
and explained it by saturation. Saturation would make R13 vanish along with R17; R13 does
not vanish. The manuscript uses R14 together with R17 as its main argument against the
threshold hypothesis; the argument should be restated, because R17 as it stands is
evidence against a level account of R13 rather than evidence about accumulation.

\begin{theorem}[R24: shape of the saturation]
\label{thm:R24}
Let $\langle\cdot\rangle_h$ denote the $h$-weighted average over trials. Then
\begin{equation}
\frac{\dd}{\dd\tau}\ln\nu=k\langle a\rangle_h,\qquad
\boxed{\;\frac{\dd^{2}}{\dd\tau^{2}}\ln\nu
=k^{2}\bigl(\Var_h(a)-1\bigr)
=\frac{\alpha^{2}}{\sigma_Y^{4}}\Var_h(\Mrv)-\frac{\alpha^{2}}{\sigma_Y^{2}}\;}
\label{eq:curv}
\end{equation}
so the delay log-rate is concave if and only if the between-trial standard deviation of
the maintained level is smaller than the within-trial maintenance noise:
$\operatorname{sd}_h(\Mrv)<\sigma_Y$. Also
$\partial\nu/\partial\mu=(a/\sigma_Y)\nu\to0$ as $\mu\to g$,
$\partial^2\nu/\partial\mu^2=\nu(a^2-1)/\sigma_Y^{2}$ (convex for $a>1$, concave for
$a<1$), and $\nu\le\Gamma\beta/2\pi$. \emph{Status: \textsc{Theorem}.}
\end{theorem}

Equation~\eqref{eq:curv} strictly generalises the previous Corollary~3, which is the
degenerate case $\Var_h=0$ and gives $-k^2$. The generalisation matters: it converts the
observed curvature into a measurement of the trial-to-trial spread of the maintained
level (Section~\ref{sec:diagnostics}), which is the quantity the note elsewhere identifies
as the model's largest deficiency.

\begin{theorem}[R15: necessary condition]
\label{thm:R15}
Because $h$ is increasing and $\Mrv\le\theta$,
$\Gamma_0F_{D,c}h(\theta+\alpha\tau)\le\nu_{D,c}\le\Gamma_0h(\theta+\alpha\tau)$, so
\begin{equation}
\nu_{D,c_H}-\nu_{D,c_L}\le\Gamma_0h(\theta+\alpha\tau)\bigl(1-F_{D,c_L}\bigr).
\end{equation}
Integrating over the delay window and writing $N_\theta=\Gamma_0\int h(\theta+\alpha\tau)\dd\tau=2.097$,
\begin{equation}
\boxed{\;1-F_{2,c_L}\;\ge\;\frac{\Delta_c N_{2}}{N_\theta}\;}
\end{equation}
With the observed $\Delta_cN_2=1.894-1.682=0.212$ this gives
$1-F_{2,0.12}\ge0.101$: \textbf{at least ten per cent of two-second, twelve-per-cent-coherence
trials must fail to reach the internal threshold.} \emph{Status: \textsc{Theorem}.}
\end{theorem}

This replaces the qualitative ``requires a substantial fraction'' of the previous
Corollary~4 with a number. It is necessary, not sufficient; the $\theta$-sweep of
Section~\ref{sec:R15} shows the sufficient requirement is considerably stronger. The
logical structure of the R15 argument should be: prove the bound, then use the accuracy
data to contradict it, and demote the $\theta$-sweep to an illustration.

\begin{theorem}[R3 and R15 are compatible only inside a narrow window]
\label{thm:R3R15}
Write $\Delta_c\mu_D$ for the coherence contrast in the level at the response signal. To
first order, premotor execution time satisfies $\partial\E[\text{PMET}]/\partial\mu=-1/w$
(Proposition~\ref{prop:pmet}), while the delay rate satisfies
$\partial\nu/\partial\mu=h'\Gamma_0$. Hence
\begin{align}
\text{R3}&\iff \Delta_c\mu_{0.3}/\Delta_c\mu_{2}>1,\\
\text{R15}&\iff \Delta_c\mu_{0.3}/\Delta_c\mu_{2}<r^{*}:=\frac{h'(\mu_2)}{h'(\mu_{0.3})}
=\frac{a_2}{a_{0.3}}\cdot\frac{\nu_2}{\nu_{0.3}} .
\end{align}
At the observed rates and distances $r^{*}=1.31$.
\emph{Status: \textsc{Theorem}, given the first-order PMET map.}
\end{theorem}

The empirical content is sharp. In J~Neuroscience the PMET coherence contrasts are
$28.2$~ms (short) against $7.1$~ms (long), a ratio of $3.97$, far outside $(1,1.31)$: R3
and R15 are incompatible in that data set, and no calibration can reconcile them. In AGPS
Experiment~2 the same contrasts are $30.5$ and $27.2$~ms, a ratio of $1.12$, which is
\emph{inside} $(1,1.31)$ --- and AGPS is the experiment in which R15 is larger and R3 is
absent. The two experiments differ precisely in whether the compression ratio exceeds
$r^{*}$. This is a considerably more informative statement than ``R3 is reproduced here
and does not replicate there''.

%----------------------------------------------------------------------
\subsection{Predictions: stimulus period (R6--R11)}
\label{sec:stimpred}

All of these follow from \eqref{eq:master-stim2}. The previous version treated them as
simulation results with one heuristic proposition; the gating lemma makes them tractable.

\begin{theorem}[R8: the stimulus-period rate ramps]
\label{thm:R8}
$\nu_{D,c}(t)$ is increasing in $t$ on the analysis window, because
(i) $\Pr(\Tstop\le t)$ is non-decreasing;
(ii) $\rho(\tau;A)$ is increasing in $\tau$ wherever $A>\sigma_V(\tau)$, since
$\partial_s[s^{-1}\varphi(A/s)]=s^{-2}\varphi(A/s)(A^2/s^2-1)$, and the covariance term
$\gamma=K(\tau)^2/2$ adds a further positive conditional drift $\gamma A/s^2$;
(iii) in the $0.3$-second condition preparation adds $\alpha(t-0.52)^{+}$ from inside the
window. \emph{Status: \textsc{Theorem}.}
\end{theorem}

\begin{proposition}[R6, R9: exact decomposition]
\label{prop:R69}
Couple the two duration conditions on the same $W_x$; they are identical up to $0.52$~s.
For $t>0.52$,
\begin{equation}
\nu_{0.3,c}(t)-\nu_{2,c}(t)=\underbrace{\E\bigl[\mathbf 1\{T_\theta\le0.52\}
\bigl(\rho(t-T_\theta;g-\theta-\alpha(t-0.52))-\rho(t-T_\theta;g-\theta)\bigr)\bigr]}_{\text{(i) preparation, }\ge0}
\end{equation}
\begin{equation}
+\underbrace{\E\bigl[\mathbf 1\{T_\theta>0.52\}\rho(t-0.52;\,g-|\tilde y(0.52)|-\alpha(t-0.52))\bigr]}_{\text{(ii) sub-threshold committers, }\ge0}
-\underbrace{\E\bigl[\mathbf 1\{0.52<T_\theta\le t\}\rho(t-T_\theta;g-\theta)\bigr]}_{\text{(iii) late threshold-reachers, }\ge0}.
\end{equation}
R6 is the statement that (i)$+$(ii)$>$(iii) integrated over the window; R9 is the same
statement for the time derivative. Term (i) grows multiplicatively at rate $k\,a$; term
(iii) grows at the rate of the commitment hazard after $0.52$~s.
\emph{Status: \textsc{Decomposition}.}
\end{proposition}

This supersedes the old Proposition~6. The mechanism is not that ``the level rises at
$\bar v$ for uncommitted trials''; uncommitted trials contribute nothing. R9 is the
competition between the preparation ramp in the short condition and the continuing
commitment hazard in the long one. It yields a directly testable consequence: R9 should be
carried by the post-$0.52$~s portion of the window and should be small or absent at high
coherence, where the commitment hazard in the 2-second condition is already exhausted by
$0.52$~s.

\begin{theorem}[R7: main effect of coherence]
\label{thm:R7}
If $T_\theta^{c_H}\stle T_\theta^{c_L}$ (A4) and $\tau\mapsto\rho(\tau;g-\theta)$ is
non-decreasing (Theorem~\ref{thm:R8}(ii)), then \eqref{eq:stim2s} gives
$\nu_{2,c_H}(t)\ge\nu_{2,c_L}(t)$ pointwise. In the $0.3$-second condition the same
argument applies to the $\theta$-reaching component and (A3) to the remainder.
\emph{Status: \textsc{Conditional} on (A3), (A4); (A4) proved in Lemma~\ref{lem:girsanov}
in the unfiltered limit.}
\end{theorem}

\begin{proposition}[R10: two components of opposite sign --- a new prediction]
\label{prop:R10}
In the $0.3$-second condition the ramp after $0.52$~s is multiplicative, so its absolute
slope is proportional to the rate, which is larger at high coherence: this produces R10
with the observed sign. In the $2$-second condition the rise is driven by $F_c(t)$, which
saturates \emph{earlier} at high coherence, producing the opposite sign.
\emph{Status: \textsc{Decomposition}, with the prediction:}
\begin{quote}
\emph{R10 should be positive when computed within the $0.3$-second condition alone, and
zero or negative when computed within the $2$-second condition alone.}
\end{quote}
The pooled effect is observed to be small ($+0.00041$) and the model's is small
($+0.00172$), consistent with near-cancellation, but the split has never been examined.
\end{proposition}

\begin{proposition}[R11]
Nothing in \eqref{eq:master-stim2} forces the duration $\times$ coherence interaction to
vanish; the model predicts a small nonzero value whose sign depends on how fast $F_{2,c}$
saturates relative to the sub-threshold spread at $0.3$~s.
\emph{Status: \textsc{Not signed}; predicted small.}
\end{proposition}

%----------------------------------------------------------------------
\subsection{Predictions: latencies and channel selectivity (R1--R5)}
\label{sec:latency}

\begin{proposition}[The PMET map]
\label{prop:pmet}
Let $\mu_R=\Mrv+\alpha\,(t_R+T_{\text{enc}}-t_a)$ be the level at the moment the second
accumulation begins. Ignoring the maintenance fluctuation, the onset of the response burst
is the first passage of $g$ by $\mu_R+q$ with $\dd q=w\dd t+\sigma_q\dd W_q$, so
\begin{equation}
\E[\text{PMET}\mid\mu_R]=T_{\text{enc}}+\frac{g-\mu_R}{w},\qquad
\frac{\partial\E[\text{PMET}]}{\partial\mu_R}=-\frac1w ,
\end{equation}
the first identity being exact for a one-sided barrier with positive drift. The
maintenance fluctuation $V$ shortens this by a term independent of $(D,c)$ to first order.
\end{proposition}

\begin{theorem}[R1, R2]
Under (A2), (A3), $\mu_R$ is stochastically larger after a longer stimulus and at higher
coherence, so premotor execution time is shorter in both.
\emph{Status: \textsc{Conditional} on (A2), (A3).}
\end{theorem}

\begin{proposition}[R3]
$\Delta_c\E[\text{PMET}]_D=-\Delta_c\E[\mu_R]_D/w$, so R3 is exactly (A5). The model does
not determine the sign of (A5); the Methods of the manuscript say so explicitly. R3 should
therefore be reported as evidence against the temporal-preparation hypothesis and not as
evidence for the combined hypothesis. \emph{Status: \textsc{Not predicted} (sign).}
\end{proposition}

\begin{proposition}[R4: motor execution time is invariant by construction, to first order]
Motor execution time is the first passage from $g$ to $g+r$ plus $T_r$. Conditionally on
being at $g$, the subsequent dynamics are driven by $q$ and $V$, whose laws contain no
$(D,c)$-dependent parameter, so $T_{\text{motor}}\perp(D,c)$ exactly at first order.
A second-order dependence survives, because the \emph{decomposition} of $g$ into
$\mu+q^{+}+V$ at onset differs across conditions: a trial arriving at $g$ with a large
mean-reverting $V$ and a small $q$ is likelier to fall back. The model therefore predicts
a small residual, which is what is observed ($2.5$~ms on $160$~ms, $1.6$~per~cent).
\emph{Status: \textsc{Structural} at first order.} Reproducing R4 in simulation is not
evidence for the model and should not be counted as such.
\end{proposition}

\begin{proposition}[R5]
The model represents only the selected channel; the non-selected channel carries
$-s\,y$, which never approaches $g$. \emph{Status: \textsc{Structural}. Not a test.}
\end{proposition}

%----------------------------------------------------------------------
\subsection{Predictions: response-signal and response periods (R18--R23)}
\label{sec:response}

The previous version had no theory here. The key observation is that the second
accumulation does not begin until $t_R+T_{\text{enc}}=0.22$~s, which is \emph{after} the
$0$--$0.2$~s response-signal window. Inside that window the process is still the delay
process. What changes is the risk set.

\begin{theorem}[R18: attrition, not dynamics]
\label{thm:R18}
Let $S(t)=\Pr(\text{PMET}>t)$ and let $\eta(t)=-\dd\ln S/\dd t$ be the hazard of response
onset. Since a trial contributes no further partial activations once its response burst
has begun, the per-trial observed rate is $\nu^{\text{obs}}(t)=\nu^{\text{free}}(t)S(t)$
and
\begin{equation}
\boxed{\;\frac{\dd}{\dd t}\ln\nu^{\text{obs}}=k\langle a\rangle_h-\eta(t)\;}
\end{equation}
R18 (a decrease) holds if and only if $\eta(t)>k\langle a\rangle_h\approx0.7$~s$^{-1}$.
With the observed PMET distribution (mean $\approx350$~ms, sd $\approx180$~ms),
$\eta$ rises from $0.35$~s$^{-1}$ at $t=0$ to $1.96$~s$^{-1}$ at $t=0.2$~s, so R18 holds
over most of the window. \emph{Status: \textsc{Theorem}, parameter-free given the observed
PMET distribution.}
\end{theorem}

Two consequences, both new and both testable in the existing data:
\begin{enumerate}
\item The decline should \emph{accelerate} across the window and should be flat or slightly
positive in the first bin, where $\eta<k\langle a\rangle_h$.
\item Dividing the observed response-signal-locked rate by the empirical survival function
$\hat S(t)$ should restore a rising profile continuous with the delay period. If it does
not, the attrition account fails.
\end{enumerate}
This also explains why the implementation error described in Section~\ref{sec:method}
inverted R18: that error changed exactly the treatment of the terminal burst, which is the
attrition term.

\begin{proposition}[R19]
At $t=0$ the ratio $\nu_{2,c}/\nu_{0.3,c}$ equals its delay-period value (continuity), and
the two survival functions differ only by the PMET difference of $50$~ms. R19 therefore
follows from R12 with a small correction. \emph{Status: \textsc{Conditional} on (A2).}
\end{proposition}

\begin{proposition}[R20]
$\dd\ln\nu^{\text{obs}}_{2}/\dd t-\dd\ln\nu^{\text{obs}}_{0.3}/\dd t=-[\eta_2(t)-\eta_1(t)]<0$
because premotor execution time is shorter after a 2-second stimulus (R1), so its hazard is
uniformly higher. \emph{Status: \textsc{Theorem}, given R1.} Note the linkage: the sign of
R20 is forced by the sign of R1. At the observed PMET means the hazard difference at
$t=0.15$~s is $0.52$~s$^{-1}$.
\end{proposition}

\begin{proposition}[R21 is a power null, not a structural zero]
The coherence contrast in this window equals the delay coherence contrast (R13) modulated
by the coherence difference in hazard (from R2), which acts in the opposite direction. The
predicted magnitude is $\approx(\text{R19}/\text{R12})\times\text{R13}\approx0.012$ per
bin, about one third of R19. \emph{Status: \textsc{Decomposition}; the model does not
predict an exact zero and the observed null is consistent with insufficient power.}
\end{proposition}

\begin{theorem}[R22: the response-locked rate must vanish at the response]
\label{thm:R22}
Align at the response onset $t=0$. A counted partial activation at $t<0$ is an upcrossing
of $g$ that is followed by a downcrossing before $0$, since the excursion that reaches
$g+r$ is the response and is not counted. Hence
\begin{equation}
\nu^{\text{resp}}(t)=\nu(t)\Pr(\text{excursion beginning at }t\text{ ends before }0)
\;\longrightarrow\;0\quad\text{as }t\uparrow0,
\end{equation}
with a decay scale equal to the mean excursion duration. R22 is therefore structural, and
its \emph{scale} is predicted: $\bar D=\Phi(-a)/\nu\approx78$~ms at calibrated parameters.
\emph{Status: \textsc{Theorem}.}
\end{theorem}

\begin{proposition}[R23]
The decline of Theorem~\ref{thm:R22} is multiplicative and common to conditions, so an
absolute-rate interaction of sign $-\operatorname{sign}(\nu_2-\nu_{0.3})$ follows
immediately from R19. \emph{Status: \textsc{Theorem}, given R19.}
\end{proposition}

%----------------------------------------------------------------------
\subsection{Excursion duration}

\begin{proposition}[Duration identity]
$\dfrac{\E A[0,T]}{\E N[0,T]}=\dfrac{\int_0^T\Phi(-a(t))\dd t}{\int_0^T\nu(t)\dd t}$.
\end{proposition}
The proof is unchanged. Two clarifications:
\begin{enumerate}
\item The ratio of expectations is not the expectation of a ratio, but it \emph{is} the
estimand of the pooled empirical estimator (total time above threshold divided by total
number of bursts, pooled over trials), which is how the quantity is measured. The previous
version's interpretation is therefore correct for the intended use, and the critique that
$\E A/\E N$ is not ``the mean duration of a randomly selected excursion'' does not bite.
\item The boundary correction is $O(1/\E N)$ relative, and $\E N\approx1$ per trial in the
delay window, so it is not asymptotically negligible: the numerical check gives $2$~per
cent. The phrase ``up to the contribution of at most one incomplete excursion at each
endpoint'' should be accompanied by that number.
\end{enumerate}

%----------------------------------------------------------------------
\subsection{Quantitative diagnostics}
\label{sec:diagnostics}

These are new. Each is a number that the theory extracts from the published summary
statistics without any simulation.

\paragraph{D1. The standardised distance to the gate.}
From $\nu=\nu_{\max}e^{-a^2/2}$ with $\nu_{\max}=2.506$~s$^{-1}$ and the observed
delay-window totals (per trial, over the 2~s window):
$a=1.74$ (0.3~s, pooled), $a=1.46$ (2~s, pooled), $a=1.40$ (2~s, $c=0.50$).
Two independent checks agree: (i) the observed log-slope of the delay rate,
$+0.809$ then $+0.702$~s$^{-1}$, gives $a=\text{slope}/k=1.93$ then $1.67$; (ii) the
premature-response rate of $1.35$~per~cent requires
$\nu(g+r)\approx0.003$~s$^{-1}$, i.e.\ $a+\rho_r\approx3.7$ with $\rho_r=2.11$, hence
$a\approx1.6$.
\textbf{The value $a\in[2.2,3.1]$ quoted in the previous version is not consistent with
the note's own rates}; at $a=2.5$ the predicted premature rate is $0.02$~per~cent against
$1.35$ observed. Consequences: the absorption correction is about $1$~per~cent rather than
below $0.2$; and $a$ straddles $\sqrt3=1.732$, which is the sign-change point of $h'''$.

\paragraph{D2. The trial-to-trial spread of the maintained level, measured from curvature.}
Equation~\eqref{eq:curv} with the observed change in log-slope, $-0.107$~s$^{-2}$, and
$k$ estimated from the slope itself gives
$\Var_h(\Mrv)/\sigma_Y^{2}=1-0.107/k^2$, i.e.\
$\operatorname{sd}_h(\Mrv)\approx0.24$ in units of $\sigma=1$, about $0.75\sigma_Y$. This
is an EMG-derived estimate of exactly the quantity the model is suspected of
underestimating. It should be compared with the simulated $\operatorname{sd}(\Mrv)$, which
is small because most trials saturate at $\theta$.

\paragraph{D3. The premotor-execution-time / rate consistency test.}
Both PMET and the delay rate are monotone functions of the same maintained level, so each
implies a value of $\Delta\mu$. Using $\partial\text{PMET}/\partial\mu=-1/w$ and
$\Delta\ln\nu=a\Delta\mu/\sigma_Y$:

\begin{center}
\begin{tabular}{lrrr}
\toprule
contrast & $\Delta\mu$ from PMET & $\Delta\mu$ from delay rate & ratio\\
\midrule
duration (pooled) & $0.308$ & $0.090$ & $3.4$\\
coherence $12\to50$, $0.3$~s & $0.174$ & $0.009$ & $19.8$\\
coherence $12\to50$, $2$~s & $0.044$ & $0.026$ & $1.7$\\
\bottomrule
\end{tabular}
\end{center}

\textbf{The two measures are quantitatively irreconcilable in the $0.3$-second condition,
by a factor of twenty.} The data require the maintained level to be nearly identical
across coherence after a short stimulus (the delay-rate contrast is $5$~per~cent) while
premotor execution time differs by $28$~ms. This is a stronger and more localised
statement than ``the model is uniformly too reactive to condition'', and it identifies
where the model must be changed.

\paragraph{D4. What D3 implies for the model.}
A quantity that shifts premotor execution time without shifting the delay-period level is
required. The natural candidate is already named in the manuscript: a condition-dependent
\emph{starting point of the second accumulation}, distinct from the maintained motor level
and added only at $t_R+T_{\text{enc}}$. Such a term
\begin{itemize}
\item shifts PMET (R1--R3) by $-\Delta q_0/w$;
\item leaves the delay-period rate untouched, hence leaves R12--R17 to be carried by the
small level differences the EMG data actually show;
\item leaves motor execution time invariant (R4), because it acts before the burst onset
and MET is measured from $g$;
\item predicts that the duration and coherence contrasts on the partial-activation rate
should \emph{step up} only after $T_{\text{enc}}=0.22$~s post-signal, i.e.\ outside the
$0$--$0.2$~s response-signal window and inside the response-locked window.
\end{itemize}
The last item is a sharp falsifiable prediction that distinguishes this variant from the
present one. It has not been tested.

\paragraph{D5. The $d'$ benchmark, computed for the actual process.}
For \eqref{eq:acc} without stopping,
$\E\tilde y(D)=v(D-(1-e^{-\lambda D})/\lambda)$ and
$\Var\tilde y(D)=\sigma^2[D-2(1-e^{-\lambda D})/\lambda+(1-e^{-2\lambda D})/(2\lambda)]$,
so the $d'$ ratios predicted by unbounded accumulation \emph{in this model} are
$1.871$ ($1000$ vs $300$~ms) and $1.422$ ($2000$ vs $1000$~ms), against $1.826$ and
$1.414$ for the unfiltered diffusion and $1.26$ and $1.01$ observed. The filter makes the
prediction slightly \emph{larger}, so the argument that accumulation terminates is
strengthened, not weakened, by using the correct process. The previous version should
either quote these numbers or label the $\sqrt{T}$ figures explicitly as a diffusion
benchmark.

\paragraph{D6. Jensen: the variance conjecture points the wrong way for $a>\sqrt3$.}
Expanding \eqref{eq:master-delay} to second order,
$\nu\approx\Gamma_0h(\bar m)[1+(a^2-1)\Var(\Mrv)/2\sigma_Y^2]$, so a condition contrast at
common variance $V$ obeys
$\Delta\nu\approx\Gamma_0\bigl[h'(\bar m)+\tfrac{V}{2}h'''(\bar m)\bigr]\Delta\bar m$.
By \eqref{eq:hderiv}, $h'''>0$ for $a>\sqrt3$ and $h'''<0$ for $a<\sqrt3$. Hence:
\begin{quote}
\emph{Where $a>\sqrt3$, larger trial-to-trial variance in the maintained level
\textbf{amplifies} condition effects; where $a<\sqrt3$ it attenuates them.}
\end{quote}
The note conjectures that the model's $45$~per~cent underestimate of premotor-execution-time
variability explains its overestimate of seven condition effects. At $a>\sqrt3$ that
conjecture is backwards: too little variance would predict effects that are too
\emph{small}. Since the calibrated $a\in[1.40,1.74]$ straddles $\sqrt3=1.732$, the sign
differs across cells, and the conjecture cannot be evaluated by intuition. If the variance
deficit does explain the overestimation, it must do so through its effect on the
\emph{means} --- for example, trial-to-trial drift variability pushes more trials to
saturation at $\theta$ in both conditions and compresses $\Delta\bar m$ --- and not through
Jensen. That is a different and separately testable mechanism.

%----------------------------------------------------------------------
\subsection{Status of R1--R24}

\begin{center}
\begin{tabular}{llp{0.46\textwidth}}
\toprule
& status & basis\\
\midrule
R1 & Conditional (A2) & Prop.~\ref{prop:pmet}, $\partial\text{PMET}/\partial\mu=-1/w$\\
R2 & Conditional (A3) & same\\
R3 & Not predicted (sign) & equals (A5); manuscript's Methods concur\\
R4 & Structural (1st order) & no $(D,c)$ parameter after $g$; residual $\approx1.6$\%\\
R5 & Structural & architecture; not a test\\
R6 & Decomposition & Prop.~\ref{prop:R69}, terms (i)+(ii) vs (iii)\\
R7 & Conditional (A3),(A4) & Thm.~\ref{thm:R7}; (A4) proved unfiltered\\
R8 & Theorem & Thm.~\ref{thm:R8}\\
R9 & Decomposition & Prop.~\ref{prop:R69}; old Prop.~6 mechanism corrected\\
R10 & Decomposition & Prop.~\ref{prop:R10}; new within-duration prediction\\
R11 & Not signed & predicted small\\
R12 & Conditional (A2) & Thm.~\ref{thm:R1213}\\
R13 & Conditional (A3) & Thm.~\ref{thm:R1213}\\
R14 & \textbf{Theorem} & Thm.~\ref{thm:R14}, unconditional\\
R15 & \textbf{Theorem} (bound) & Thm.~\ref{thm:R15}: $1-F_{2,0.12}\ge0.101$; Thm.~\ref{thm:R3R15}\\
R16 & Conditional (A2) & Thm.~\ref{thm:R16}, needs only $a>1$\\
R17 & \textbf{Theorem} (against) & Cor.~\ref{cor:R17}: incompatible with R13\\
R18 & \textbf{Theorem} & Thm.~\ref{thm:R18}, from observed PMET hazard\\
R19 & Conditional (A2) & continuity with R12\\
R20 & \textbf{Theorem} given R1 & hazard ordering\\
R21 & Decomposition & predicted $\approx0.012$/bin; power null\\
R22 & \textbf{Theorem} & Thm.~\ref{thm:R22}; scale $=\bar D\approx78$~ms\\
R23 & \textbf{Theorem} given R19 & multiplicative decline\\
R24 & \textbf{Theorem} & Thm.~\ref{thm:R24}, incl.\ curvature identity\\
\bottomrule
\end{tabular}
\end{center}

Nine results are now theorems of the model or constraints on the data, against one
(R24) previously. Four more are conditional on two checkable properties of the
accumulation stage. Three are decompositions that generate new tests. The headline
``sixteen of nineteen directional effects come out'' understates what the model delivers
and overstates what the simulation establishes.

%----------------------------------------------------------------------
\subsection{Errata relative to the previous Section~4}
\label{sec:errata}

\begin{enumerate}
\item \textbf{H1.} The general Gaussian regularity theorem invoked is false as stated.
Replaced by the pathwise argument (Prop.~\ref{prop:C1}).
\item \textbf{$C^2$ but not $C^4$} $\to$ $C\in C^2(\mathbb R)\setminus C^3(\mathbb R)$,
because $C'''$ is odd and $C'''(0^{+})\neq0$.
\item \textbf{$sU\stackrel{d}{=}U$.} Justified by independence (A1), not by symmetry of
$\varepsilon$; $s$ is measurable with respect to the accumulation filtration.
\item \textbf{Transient bound.} The relative variance deficit is $6.056\,e^{-2\kappa\tau}$,
not $e^{-2\kappa\tau}$. Also $\Cov(V,\dot V)(\tau)=K(\tau)^2/2$ exactly, which is what the
reduced formula requires to vanish.
\item \textbf{Absorption.} $\nu^{\text{obs}}\le\nu$ is asserted without proof and the two
conditioning effects have opposite signs, since $\Cov(V(u),\dot V(t))<0$ for $u<t$.
Replaced by the proved bound $\nu^{\text{obs}}\le\nu/(1-p)$.
\item \textbf{The value of $a$.} $[2.2,3.1]$ is inconsistent with the note's own rates and
premature-response rate; the correct range is $[1.40,1.74]$. The absorption correction is
therefore about $1$~per~cent, not below $0.2$.
\item \textbf{Upper bound on the rate.} $\Gamma\beta/2\pi=2.506$, not $\beta/2\pi=2.421$.
\item \textbf{Discretisation.} The Euler bias is $\delta=\lambda\kappa\Delta/(\lambda+\kappa)$
on $\sigma_Y^2$ and enters $\ln\nu$ as $(a^2-1)\delta/2$; it is therefore \emph{not}
condition-independent. Use exact transition sampling.
\item \textbf{Proposition~6 (R9).} The mechanism is misattributed: pre-commitment trials
contribute exactly zero crossings because $\theta<g$. Replaced by Prop.~\ref{prop:R69}.
\item \textbf{Duration identity.} The ratio-of-expectations reading is correct for the
pooled estimator, but the boundary correction is $O(1/\E N)$ with $\E N\approx1$ and should
be quoted as the measured $2$~per~cent.
\item \textbf{$d'$ benchmark.} Computed for the filtered process the ratios are $1.871$
and $1.422$, not $1.826$ and $1.414$.
\item \textbf{Corollary on R15.} ``Requires a substantial fraction'' replaced by
$1-F_{2,c_L}\ge\Delta_cN_2/N_\theta=0.101$.
\item \textbf{The variance conjecture.} At $a>\sqrt3$ the Jensen term implies that too
little variance produces effects that are too \emph{small}, the opposite of the note's
conjecture. The calibrated $a$ straddles $\sqrt3$; the conjecture must be evaluated
numerically, not asserted.
\item \textbf{R17.} Not a corollary of R15 explained by saturation. Given R13 and the
proportionality law \eqref{eq:proportionality}, R17 is incompatible with any level-based
account of the coherence effect.
\end{enumerate}

%----------------------------------------------------------------------
\subsection{Recomputed constants}
\label{sec:numbers}

\begin{center}
\begin{tabular}{lrl}
\toprule
quantity & value & source\\
\midrule
$\sigma_Y$ & $0.32014$ & $\sqrt{\lambda\sigma_m^2/(2\kappa(\lambda+\kappa))}$\\
$\sigma_{\dot Y}$ & $4.86988$ & $\sqrt{\lambda^2\sigma_m^2/(2(\lambda+\kappa))}$\\
$\beta=\sqrt{\lambda\kappa}$ & $15.212$~s$^{-1}$ & \\
$\Gamma$ & $1.03497$ & trend factor\\
$\Gamma_0$ & $2.01074$ & $\alpha\Phi(m)+\sigma_{\dot Y}\varphi(m)$\\
$\nu_{\max}=\Gamma\beta/2\pi$ & $2.5057$~s$^{-1}$ & \\
$k=\alpha/\sigma_Y$ & $0.41982$~s$^{-1}$ & \\
$\rho_r=r/\sigma_Y$ & $2.1088$ & \\
$(g-\theta)/\sigma_Y$ & $2.1760$ & $a$ at commitment, saturated trial\\
$a$ at delay onset / offset (saturated) & $1.756$ / $0.916$ & $\mu=\theta+\alpha\tau$, $\tau\in[1,3]$\\
$a$ inferred, $0.3$~s / $2$~s / max cell & $1.742$ / $1.462$ / $1.395$ & from observed totals\\
$N_\theta$ (saturated window total) & $2.097$ & $\Gamma_0\int_1^3h(\theta+\alpha\tau)\dd\tau$\\
transient prefactor & $6.0564$ & $\lambda(\lambda+\kappa)/(\lambda-\kappa)^2$\\
$\delta/\Delta$ (Euler) & $7.1758$~s$^{-1}$ & $\lambda\kappa/(\lambda+\kappa)$\\
$r^{*}$ (R3/R15 window) & $1.314$ & $(a_2/a_{0.3})(\nu_2/\nu_{0.3})$\\
$d'$ ratios, filtered & $1.871$, $1.422$ & $1000/300$, $2000/1000$~ms\\
\bottomrule
\end{tabular}
\end{center}

%----------------------------------------------------------------------
\subsection{Open questions the theory does not close}

\begin{enumerate}
\item (A3) and (A5) are assumed, not proved. Lemma~\ref{lem:coupling} gives the pathwise
ordering of the \emph{signed} $m_0$; the passage to $|m_0|$ fails only on sub-threshold
error trials, and the residual is bounded by $\sup h\cdot\Pr(\Mrv<\theta,\,m_0<0)$. This is
a five-line simulation check and should be done.
\item The observed R18 magnitude has not been reconstructed from the observed PMET
distribution. Theorem~\ref{thm:R18} makes that a one-line prediction with no free
parameter.
\item The response-stage observable is defined by the burst-detection rule, not by the
process. Nothing in Section~\ref{sec:response} survives a change in how the terminal burst
is treated, which is exactly what previously inverted R18. The rule should be stated in
the note as part of the model.
\item D3 and D4 identify a discrepancy the present model cannot absorb. Whether a
condition-dependent second-accumulation starting point resolves it, without breaking R4 or
R19--R23, has not been tested.
\end{enumerate}

% %======================================================================
% \section{The model}
% \label{sec:model}

% \subsection{Variables and phases}

% \begin{itemize}
% \item $x(t)$ accumulated motion evidence; $\tilde y(t)$ the filtered motor trace; $y(t)$
% total motor preparatory activity.
% \item $\lambda$ filter rate, fixed at 21.47 per second; $\sigma$ accumulation noise
% coefficient, fixed at 1, which sets the scale.
% \item $v_c = A(c/0.5)^{0.6}$ drift rate at coherence $c$.
% \item $\theta$ internal threshold, applied to $\tilde y$; $g$ gating inhibition; $r$ the
% further amount above $g$ needed to press; $T_r$ a residual delay.
% \item $\kappa,\sigma_m$ relaxation rate and amplitude of the delay-period perturbation.
% \item $\alpha$ the slope of motor preparation; $w,\sigma_q$ drift and noise of the second
% accumulation; $T_{\text{enc}}=0.22$ s the encoding time, common to stimulus and signal.
% \item $z(t)=s\,y(t)$ activity in the selected channel, $s=\pm1$.
% \end{itemize}

% Encoded evidence is available on $[T_{\text{enc}},T_{\text{enc}}+D]$. Before the internal
% threshold is reached,
% \begin{equation}
% \dd x = v_c\dd t + \sigma\dd W_x,\qquad \dot{\tilde y} = \lambda(x-\tilde y).
% \end{equation}
% Accumulation stops at
% $\Tstop=\min\bigl(\inf\{t:|\tilde y|\ge\theta\},\,T_{\text{enc}}+D\bigr)$, with
% $s=\operatorname{sign}\tilde y(\Tstop)$ and $m_0=\tilde y(\Tstop)$. From $\Tstop$ the filter
% input becomes $m_0+\varepsilon$ with
% $\dd\varepsilon=-\kappa\varepsilon\dd t+\sigma_m\dd W_\varepsilon$. From
% $t_R+T_{\text{enc}}$ a second accumulation $\dd q=w\dd t+\sigma_q\dd W_q$ adds $s\,q^{+}$.
% The button is pressed when $z$ reaches $g+r$, plus $T_r$.

% \subsection{Motor preparation is added outside the filter}
% \label{sec:outside}

% \begin{equation}
% y(t)=\tilde y(t)+s\,\alpha\,\bigl(\min(t,\,t_R+T_{\text{enc}})-t_{\text{anchor}}\bigr)^{+}.
% \end{equation}

% \begin{remark}[Why outside]
% The filter exists to recover a decision variable corrupted by transmission noise. Temporal
% preparation is not a noisy measurement of anything, so there is no reason for it to be
% filtered. Empirically the change is neutral: with everything else held, the fit moves in the
% third decimal place and one further effect becomes stable, sixteen of nineteen rather than
% fifteen. Theoretically it is a simplification: $\sigma_Y$ and $\beta$ are then exactly those
% of the delay-period perturbation alone, with no contamination by the ramp, so every result
% in Section~\ref{sec:theory} holds without approximation.
% \end{remark}

% \subsection{The two anchorings}

% \begin{description}
% \item[Offset anchoring] $t_{\text{anchor}}=T_{\text{enc}}+D$.
% \item[Commitment anchoring] $t_{\text{anchor}}=\Tstop$, as in the original statement of the
% third hypothesis.
% \end{description}

% Commitment anchoring is the more natural of the two on general grounds: a system that has
% committed has no reason to wait for an external event before preparing, and with a long
% stimulus and a short delay, deferring preparation to the offset would be implausible.
% Section~\ref{sec:anchor} reports that the data nonetheless require offset anchoring in this
% design, and explains why.

% Under offset anchoring, a trial reaching $\theta$ before the stimulus ends sits at a
% \emph{plateau}: accumulation has stopped, preparation has not started, the level is constant
% at $m_0$ plus noise. At calibrated parameters this lasts 0.04 s after a 0.3-second stimulus
% and 1.1 to 1.7 s after a 2-second stimulus. It is what produces R9.

% %======================================================================
% \section{Theory}
% \label{sec:theory}

% This section is written to be checkable line by line. Section~\ref{sec:setup} fixes the
% process, Section~\ref{sec:cov} computes its covariance in closed form,
% Section~\ref{sec:KR} states the Kac--Rice theorem and verifies its hypotheses,
% Sections~\ref{sec:apply}--\ref{sec:absorb} apply it and quantify the two corrections that
% the application requires, and Section~\ref{sec:durid} derives the duration identity without
% any appeal to ergodicity. Every quantitative claim is checked against direct simulation at
% the calibrated parameters, reported inline.

% Throughout, the calibrated values are $\lambda=21.47$, $\kappa=10.778$, $\sigma_m=1.8216$,
% $\alpha=0.1344$, $g=1.5897$, $r=0.6751$, $\theta=0.8931$, all in the units fixed by
% $\sigma=1$.

% \subsection{The delay-period process}
% \label{sec:setup}

% Fix a trial that has committed at $\Tstop$ with $s=\operatorname{sign}\tilde y(\Tstop)$ and
% $m_0=\tilde y(\Tstop)$. For $t>\Tstop$ the filter input is $m_0+\varepsilon(t)$ with
% $\varepsilon(\Tstop)=0$ and
% $\dd\varepsilon=-\kappa\varepsilon\dd t+\sigma_m\dd W$, and preparation is added outside
% the filter. Activity in the selected channel is therefore
% \begin{equation}
% z(t)=s\,\tilde y(t)+\alpha\,(t-t_{\text{off}})^{+}.
% \end{equation}

% Write $\tilde y=m_0+U$ where $U$ is the filter output driven by $\varepsilon$ alone, so that
% $s\tilde y=|m_0|+sU$. Because $\varepsilon$ is a symmetric process and the filter is linear,
% $sU$ has the same law as $U$. Hence
% \begin{equation}
% z(t)=\mu(t)+V(t),\qquad
% \mu(t)=|m_0|+\alpha\,(t-t_{\text{off}})^{+},\qquad
% V \stackrel{d}{=} U,
% \label{eq:decomp}
% \end{equation}
% with $V$ a centred Gaussian process and $\mu$ deterministic given $m_0$. This is the only
% place where placing preparation outside the filter matters: it puts the entire effect of
% preparation into $\mu$ and leaves the law of $V$ untouched.

% \subsection{Covariance of $V$ in closed form}
% \label{sec:cov}

% \begin{lemma}[Kernel]
% \label{lem:kernel}
% For $\tau=t-\Tstop\ge0$,
% \begin{equation}
% V(\Tstop+\tau)=\int_0^\tau K(\tau-p)\dd W(p),\qquad
% K(v)=\frac{\lambda\sigma_m}{\lambda-\kappa}\bigl(e^{-\kappa v}-e^{-\lambda v}\bigr).
% \end{equation}
% \end{lemma}

% \begin{proof}
% Solving the Ornstein--Uhlenbeck equation from $\varepsilon(\Tstop)=0$ gives
% $\varepsilon(\Tstop+\rho)=\sigma_m\int_0^{\rho}e^{-\kappa(\rho-p)}\dd W(p)$. The filter
% $\dot U=\lambda(\varepsilon-U)$ with $U(\Tstop)=0$ has solution
% $U(\Tstop+\tau)=\lambda\int_0^{\tau}e^{-\lambda(\tau-\rho)}\varepsilon(\Tstop+\rho)\dd\rho$.
% Substituting and exchanging the order of integration --- legitimate by the stochastic Fubini
% theorem, since the kernel $(\rho,p)\mapsto e^{-\lambda(\tau-\rho)}e^{-\kappa(\rho-p)}
% \mathbf 1\{p\le\rho\le\tau\}$ is bounded on a bounded domain --- gives
% \begin{equation}
% U(\Tstop+\tau)=\lambda\sigma_m\int_0^{\tau}
% \Bigl(\int_p^{\tau}e^{-\lambda(\tau-\rho)}e^{-\kappa(\rho-p)}\dd\rho\Bigr)\dd W(p).
% \end{equation}
% The inner integral equals
% $e^{-\lambda\tau+\kappa p}\int_p^{\tau}e^{(\lambda-\kappa)\rho}\dd\rho
% =\bigl[e^{-\kappa(\tau-p)}-e^{-\lambda(\tau-p)}\bigr]/(\lambda-\kappa)$ for
% $\lambda\neq\kappa$.
% \end{proof}

% \begin{proposition}[Stationary covariance]
% \label{prop:cov}
% As $\tau\to\infty$ the law of $V$ converges to a stationary centred Gaussian process with
% \begin{equation}
% C(\tau)=\Cov\bigl(V(t),V(t+\tau)\bigr)
% =B\left[\frac{e^{-\kappa|\tau|}}{\kappa}-\frac{e^{-\lambda|\tau|}}{\lambda}\right],
% \qquad B=\frac{\lambda^{2}\sigma_m^{2}}{2(\lambda-\kappa)(\lambda+\kappa)} .
% \label{eq:C}
% \end{equation}
% Consequently
% \begin{equation}
% \sigma_Y^{2}:=C(0)=\frac{\lambda\sigma_m^{2}}{2\kappa(\lambda+\kappa)},\qquad
% C'(0^{+})=0,\qquad
% \sigma_{\dot Y}^{2}:=-C''(0)=\frac{\lambda^{2}\sigma_m^{2}}{2(\lambda+\kappa)},
% \end{equation}
% and $\beta:=\sigma_{\dot Y}/\sigma_Y=\sqrt{\lambda\kappa}$.
% \end{proposition}

% \begin{proof}
% By Lemma~\ref{lem:kernel} and the Itô isometry, for $\tau\ge0$,
% $C(\tau)=\lim_{t\to\infty}\int_0^{t}K(v)K(v+\tau)\dd v=\int_0^{\infty}K(v)K(v+\tau)\dd v$.
% Expanding the product,
% \begin{equation}
% K(v)K(v+\tau)=\frac{\lambda^2\sigma_m^2}{(\lambda-\kappa)^2}
% \Bigl[e^{-\kappa\tau}e^{-2\kappa v}-e^{-\lambda\tau}e^{-(\lambda+\kappa)v}
% -e^{-\kappa\tau}e^{-(\lambda+\kappa)v}+e^{-\lambda\tau}e^{-2\lambda v}\Bigr].
% \end{equation}
% Integrating term by term over $v\in[0,\infty)$,
% \begin{equation}
% \int_0^{\infty}\!K(v)K(v+\tau)\dd v
% =\frac{\lambda^2\sigma_m^2}{(\lambda-\kappa)^2}\left[
% e^{-\kappa\tau}\Bigl(\frac{1}{2\kappa}-\frac{1}{\lambda+\kappa}\Bigr)
% +e^{-\lambda\tau}\Bigl(\frac{1}{2\lambda}-\frac{1}{\lambda+\kappa}\Bigr)\right].
% \end{equation}
% Now $\frac{1}{2\kappa}-\frac{1}{\lambda+\kappa}=\frac{\lambda-\kappa}{2\kappa(\lambda+\kappa)}$
% and $\frac{1}{2\lambda}-\frac{1}{\lambda+\kappa}=-\frac{\lambda-\kappa}{2\lambda(\lambda+\kappa)}$,
% which gives \eqref{eq:C} after collecting $B$. Evenness in $\tau$ follows from stationarity.

% For the derivatives, on $\tau>0$ we have
% $C'(\tau)=B[-e^{-\kappa\tau}+e^{-\lambda\tau}]$ and
% $C''(\tau)=B[\kappa e^{-\kappa\tau}-\lambda e^{-\lambda\tau}]$. Letting $\tau\downarrow0$
% gives $C'(0^{+})=B(-1+1)=0$ and $C''(0)=B(\kappa-\lambda)=-B(\lambda-\kappa)$, that is
% $-C''(0)=\lambda^2\sigma_m^2/(2(\lambda+\kappa))$. Setting $\tau=0$ in \eqref{eq:C} gives
% $C(0)=B(\lambda-\kappa)/(\kappa\lambda)=\lambda\sigma_m^2/(2\kappa(\lambda+\kappa))$.
% Finally $\beta^2=-C''(0)/C(0)=\lambda\kappa$.
% \end{proof}

% \begin{remark}[$C'(0)=0$ and its meaning]
% Because $C$ is even and differentiable at the origin with $C'(0^{+})=0$, we have
% $\Cov(V(t),\dot V(t))=C'(0)=0$ \emph{exactly}, not merely asymptotically. This is what makes
% the reduced Rice formula available. Note that $C$ is $C^2$ but not $C^4$ at the origin,
% since $C'''(0^{+})=B(\lambda^2-\kappa^2)\neq0$; the process is therefore once but not twice
% differentiable, which is all that is required below.
% \end{remark}

% \paragraph{Numerical check.} At calibrated parameters, $C(0)=0.102488$ against the closed
% form $\sigma_Y^2=0.102488$; $-C''(0)=23.7158$ against $\sigma_{\dot Y}^2=23.7158$. Direct
% simulation with $6\times10^{4}$ paths gives $\Var(V)=0.10243$ and empirical autocovariances
% $0.09842$, $0.08458$, $0.05803$, $0.02260$ at lags $0.02$, $0.05$, $0.10$, $0.20$ s against
% $0.09865$, $0.08475$, $0.05797$, $0.02243$ from \eqref{eq:C}.

% \subsection{Conditions for stationarity}
% \label{sec:stat}

% The exact transient moments follow from Lemma~\ref{lem:kernel} by integrating
% $K^2$ over $[0,\tau]$ rather than $[0,\infty)$:
% \begin{equation}
% \Var V(\Tstop+\tau)=\frac{\lambda^{2}\sigma_m^{2}}{(\lambda-\kappa)^{2}}
% \left[\frac{1-e^{-2\kappa\tau}}{2\kappa}
% -\frac{2\bigl(1-e^{-(\lambda+\kappa)\tau}\bigr)}{\lambda+\kappa}
% +\frac{1-e^{-2\lambda\tau}}{2\lambda}\right].
% \end{equation}
% The three exponentials decay at rates $2\kappa$, $\lambda+\kappa$ and $2\lambda$, of which
% the slowest is $2\kappa=21.6$ per second, a time constant of $46$ ms.

% \begin{remark}[When the stationary results may be used]
% The delay window begins $1$ s after stimulus offset, and $\Tstop\le T_{\text{enc}}+D$, so
% the elapsed time $\tau$ at the start of the window is at least $1$ s, that is at least
% $21$ time constants. The relative departure from the stationary variance is then below
% $e^{-2\kappa}\approx4\times10^{-10}$. Every use of Proposition~\ref{prop:cov} below is
% inside this window and is therefore exact for practical purposes. This is the precise sense
% in which ``at stationarity'' is meant, and it does \emph{not} hold in the stimulus-period
% window, where the stationary formulas are not used.
% \end{remark}

% \subsection{The Kac--Rice theorem}
% \label{sec:KR}

% \begin{proposition}[Kac--Rice, Gaussian case]
% \label{thm:KR}
% Let $\{X(t)\}_{t\in[0,T]}$ be a Gaussian process such that
% \begin{enumerate}
% \item[(H1)] $X$ has almost surely continuously differentiable sample paths;
% \item[(H2)] for every $t$, the pair $(X(t),\dot X(t))$ has a nondegenerate Gaussian
% distribution, in particular $\Var X(t)>0$ and $\Var\dot X(t)>0$;
% \item[(H3)] $t\mapsto(\E X(t),\E\dot X(t),\Var X(t),\Var\dot X(t),\Cov(X,\dot X)(t))$ is
% continuous.
% \end{enumerate}
% Then the expected number of upcrossings of a level $u$ satisfies
% \begin{equation}
% \E\bigl[N_u[0,T]\bigr]=\int_0^{T}\!\!\int_0^{\infty}\dot x\,
% p_{X(t),\dot X(t)}(u,\dot x)\,\dd\dot x\,\dd t
% =\int_0^{T}\nu(t)\,\dd t ,
% \label{eq:KR}
% \end{equation}
% where $p_{X(t),\dot X(t)}$ is the joint density.
% \end{proposition}

% This is the standard Rice formula for Gaussian processes; see Cramér \& Leadbetter (1967),
% Chapter~10, or Azaïs \& Wschebor (2009), Theorem~3.4. It is a theorem about the
% \emph{free} process on $[0,T]$. Nothing in it refers to absorption, and nothing in it
% requires stationarity.

% \begin{lemma}[Our process satisfies H1--H3]
% \label{lem:H}
% Let $X=z$ on the delay window. Then H1--H3 hold.
% \end{lemma}

% \begin{proof}
% By \eqref{eq:decomp}, $z=\mu+V$ with $\mu$ affine, hence $C^{\infty}$, and $V$ centred
% Gaussian with covariance \eqref{eq:C}. For H1: a stationary Gaussian process has almost
% surely continuously differentiable paths whenever its covariance is twice differentiable at
% the origin with $C''$ continuous, which holds here since
% $C''(\tau)=B[\kappa e^{-\kappa|\tau|}-\lambda e^{-\lambda|\tau|}]$ is continuous on
% $\mathbb R$. Adding the deterministic $C^\infty$ function $\mu$ preserves this. For H2:
% $\Var z(t)=\sigma_Y^2>0$ and $\Var\dot z(t)=\sigma_{\dot Y}^2>0$, and since
% $\Cov(z,\dot z)=C'(0)=0$ the covariance matrix is diagonal with positive entries, hence
% nondegenerate. For H3: all five functions are constant except $\E z(t)=\mu(t)$ and
% $\E\dot z(t)=\alpha$, which are continuous.
% \end{proof}

% \subsection{Applying the theorem: the trend cannot be dropped}
% \label{sec:apply}

% \begin{proposition}[Rate with a linear trend]
% \label{prop:rate}
% Under Lemma~\ref{lem:H}, with $a(t)=(g-\mu(t))/\sigma_Y$ and $\mu_V=\alpha$,
% \begin{equation}
% \nu(t)=\frac{\varphi\bigl(a(t)\bigr)}{\sigma_Y}\,
% \Bigl[\alpha\,\Phi\!\bigl(\alpha/\sigma_{\dot Y}\bigr)
% +\sigma_{\dot Y}\,\varphi\!\bigl(\alpha/\sigma_{\dot Y}\bigr)\Bigr].
% \label{eq:ratetrend}
% \end{equation}
% Setting $\alpha=0$ recovers $\nu=\frac{\beta}{2\pi}e^{-a^2/2}$.
% \end{proposition}

% \begin{proof}
% By \eqref{eq:KR}, $\nu(t)=p_{z(t)}(g)\,\E[\dot z(t)^{+}\mid z(t)=g]$. Since
% $\Cov(z,\dot z)=0$ and the pair is jointly Gaussian, $\dot z(t)$ is independent of $z(t)$,
% so the conditioning is vacuous. The density term is
% $p_{z(t)}(g)=\sigma_Y^{-1}\varphi(a(t))$. For the slope, $\dot z(t)=\alpha+\dot V(t)$ with
% $\dot V\sim\mathcal N(0,\sigma_{\dot Y}^{2})$. Writing $\dot z=\alpha+\sigma_{\dot Y}Z$ with
% $Z$ standard normal and $m=\alpha/\sigma_{\dot Y}$,
% \begin{equation}
% \E[\dot z^{+}]=\int_{-m}^{\infty}(\alpha+\sigma_{\dot Y}\zeta)\varphi(\zeta)\dd\zeta
% =\alpha\Phi(m)+\sigma_{\dot Y}\!\int_{-m}^{\infty}\!\zeta\varphi(\zeta)\dd\zeta
% =\alpha\Phi(m)+\sigma_{\dot Y}\varphi(m),
% \end{equation}
% using $\int_{-m}^{\infty}\varphi=\Phi(m)$ and, from $\varphi'(\zeta)=-\zeta\varphi(\zeta)$,
% $\int_{b}^{\infty}\zeta\varphi(\zeta)\dd\zeta=\varphi(b)$ with $b=-m$ and $\varphi$ even.
% Setting $\alpha=0$ gives $\sigma_{\dot Y}\varphi(0)=\sigma_{\dot Y}/\sqrt{2\pi}$, and
% multiplying by $\sigma_Y^{-1}\varphi(a)$ gives the reduced form.
% \end{proof}

% \begin{corollary}[The trend acts as a level-independent multiplicative factor]
% \label{cor:trendfac}
% The ratio of \eqref{eq:ratetrend} to the $\alpha=0$ formula is
% \begin{equation}
% \Gamma=\sqrt{2\pi}\Bigl[m\Phi(m)+\varphi(m)\Bigr],\qquad m=\alpha/\sigma_{\dot Y},
% \end{equation}
% which does not depend on $a$.
% \end{corollary}

% \paragraph{Size of the correction.} At calibrated parameters $\sigma_{\dot Y}=4.870$ and
% $m=0.0276$, giving $\Gamma=1.0350$. The trend inflates every crossing rate by
% $3.5$ per cent uniformly. Simulation confirms the level-independence: the ratio of the
% trend formula to the $\alpha=0$ formula is $1.0350$ at $a=1.0$, $1.5$, $2.0$ and $2.5$
% alike.

% \begin{remark}
% Dropping the trend is therefore not harmless. It is a uniform $3.5$ per cent error at these
% parameters, and since $\Gamma$ depends only on $\alpha/\sigma_{\dot Y}$ it cancels from every
% \emph{ratio} of rates --- which is why the two-level ratio of Proposition~\ref{prop:dur} and
% all the qualitative orderings are unaffected. It does not cancel from absolute rates.
% \end{remark}

% \subsection{Absorption: the observed process is killed at $g+r$}
% \label{sec:absorb}

% Proposition~\ref{thm:KR} concerns the free process. The measured quantity is different: a
% trial that reaches $g+r$ before the response signal produces a premature response, is removed
% from the experiment, and contributes no further partial activations. The observed rate is
% therefore that of the process killed at $g+r$, per surviving trial:
% \begin{equation}
% \nu^{\text{obs}}(t)=\lim_{h\downarrow0}\frac{1}{h}\,
% \frac{\E\bigl[N_g(t,t+h)\,\mathbf 1\{\tau_{\text{abs}}>t\}\bigr]}{\Pr(\tau_{\text{abs}}>t)},
% \qquad \tau_{\text{abs}}=\inf\{t: z(t)\ge g+r\}.
% \end{equation}

% No closed form is available. Two facts bound the discrepancy. First, killing removes
% precisely the paths that have been highest, so conditioning on survival biases $z(t)$
% downward and $\nu^{\text{obs}}\le\nu$. Second, the bias is controlled by the absorption
% probability, which is itself given by the two-level ratio of Proposition~\ref{prop:dur}.

% \paragraph{Size of the correction, measured.} Simulating the delay process with and without
% absorption at $g+r$, $6\times10^{4}$ paths, $\Delta t=1$ ms:

% \begin{center}
% \begin{tabular}{lrrr}
% \toprule
% $a$ at the start of the window & absorbed & $\nu^{\text{obs}}/\nu^{\text{free}}$ & error\\
% \midrule
% 1.5 & 3.73\% & 0.9897 & $1.0$\%\\
% 2.0 & 0.69\% & 0.9981 & $0.2$\%\\
% 2.5 & 0.10\% & 0.9998 & $0.02$\%\\
% \bottomrule
% \end{tabular}
% \end{center}

% At the calibrated $a$, which lies between $2.2$ and $3.1$ across the delay, the absorption
% correction is below $0.2$ per cent and is negligible relative to the trend correction. This
% is a consequence of the observed premature-response rate being $1.35$ per cent, not an
% assumption.

% \subsection{Verification of the applied formula}

% Combining Sections~\ref{sec:apply} and~\ref{sec:absorb}, the prediction for the observed
% rate is \eqref{eq:ratetrend} with a downward correction below one per cent. Direct
% simulation of the delay process, $2.5\times10^{4}$ paths, mean level $\mu(t)=m_0+\alpha t$
% with $a(0)=2$:

% \begin{center}
% \begin{tabular}{rrrr}
% \toprule
% $\Delta t$ (ms) & empirical rate & formula \eqref{eq:ratetrend} & ratio\\
% \midrule
% 2.00 & 0.6490 & 0.6239 & 1.0402\\
% 1.00 & 0.6452 & 0.6239 & 1.0341\\
% 0.50 & 0.6430 & 0.6239 & 1.0306\\
% 0.25 & 0.6359 & 0.6239 & 1.0192\\
% \bottomrule
% \end{tabular}
% \end{center}

% The residual falls monotonically as the step shrinks, so it is discretisation, not a defect
% of the formula: an Euler scheme overstates $\Var\dot z$ because it forms a finite difference
% rather than a derivative, and the crossing count inherits the error. All simulations
% reported elsewhere in this document use $\Delta t=2$ ms, so absolute rates carry a bias of
% about $4$ per cent from this source. Every effect in Table~\ref{tab:eff} is a difference or
% a slope, in which a multiplicative bias of this kind largely cancels.

% \subsection{Excursion duration, without ergodicity}
% \label{sec:durid}

% The mean level $\mu(t)$ increases during the delay, so $z$ is \emph{not} stationary and no
% ergodic theorem applies. The identity below therefore replaces the ergodic argument used in
% earlier versions of this note, which was invalid.

% \begin{proposition}[Duration identity]
% \label{prop:dur}
% Let $A[0,T]=\int_0^T\mathbf 1\{z(t)>g\}\dd t$ be the total time above the gate and
% $N[0,T]$ the number of upcrossings. Then
% \begin{equation}
% \frac{\E\,A[0,T]}{\E\,N[0,T]}
% =\frac{\int_0^{T}\Phi\bigl(-a(t)\bigr)\dd t}{\int_0^{T}\nu(t)\dd t},
% \label{eq:durid}
% \end{equation}
% and the left-hand side is the mean excursion duration up to the contribution of at most one
% incomplete excursion at each endpoint.
% \end{proposition}

% \begin{proof}
% For the numerator, $\mathbf 1\{z(t)>g\}$ is jointly measurable and bounded, so Tonelli
% applies and $\E A=\int_0^T\Pr(z(t)>g)\dd t$. Since $z(t)\sim\mathcal N(\mu(t),\sigma_Y^2)$,
% $\Pr(z(t)>g)=\Phi(-a(t))$. For the denominator, $\E N=\int_0^T\nu(t)\dd t$ by
% Proposition~\ref{thm:KR}, whose hypotheses hold by Lemma~\ref{lem:H}. Excursions above $g$
% alternate with intervals below, so the number of complete excursions differs from $N$ by at
% most one, and the total time above differs from the sum of complete excursion durations by
% at most the two boundary excursions.
% \end{proof}

% \begin{corollary}
% \label{cor:durlocal}
% If $a$ varies slowly enough that it may be treated as constant over one excursion, the local
% mean duration is $\bar D(t)=\Phi(-a(t))/\nu(t)$.
% \end{corollary}

% \begin{remark}[When the local form is legitimate]
% The approximation in Corollary~\ref{cor:durlocal} requires the mean level to move little over
% one excursion, that is $\alpha\bar D\ll\sigma_Y$. At calibrated parameters
% $\alpha\bar D\approx0.134\times0.09=0.012$ against $\sigma_Y=0.320$, a ratio of $0.038$. The
% exact identity \eqref{eq:durid} carries no such requirement and should be preferred whenever
% the whole window is used.
% \end{remark}

% \paragraph{Numerical check.} With $\Delta t=0.25$ ms, $2.5\times10^{4}$ paths, $a(0)=2$ and
% $T=1.5$ s: empirical time above $0.07435$ s per trial against $\int\Phi(-a)\dd t=0.07300$
% (ratio $1.019$); empirical upcrossings $0.9723$ per trial against $\int\nu\,\dd t=0.9359$
% (ratio $1.039$); empirical mean duration $76.5$ ms against the predicted $78.0$ ms
% (ratio $0.980$). The two discrepancies have the same discretisation origin and partly cancel
% in the ratio.

% \subsection{The two-level ratio}

% \begin{proposition}
% \label{prop:two}
% With $\rho_r=r/\sigma_Y$, at any time $t$,
% \begin{equation}
% \frac{\nu(t;g+r)}{\nu(t;g)}
% =\exp\!\left[-\rho_r\Bigl(a(t)+\frac{\rho_r}{2}\Bigr)\right].
% \end{equation}
% \end{proposition}

% \begin{proof}
% By Proposition~\ref{prop:rate} the two rates differ only through the factor
% $\varphi(\cdot)$, since the bracket depends on $\alpha$ and $\sigma_{\dot Y}$ but not on the
% level. The standardised distance to the upper level is $(g+r-\mu)/\sigma_Y=a+\rho_r$, so the
% ratio is $\varphi(a+\rho_r)/\varphi(a)=\exp[-((a+\rho_r)^2-a^2)/2]$.
% \end{proof}

% \begin{remark}
% The trend factor $\Gamma$ of Corollary~\ref{cor:trendfac} cancels exactly here, which is why
% this proposition is unaffected by the correction of Section~\ref{sec:apply}. The same is true
% of every ratio of rates at the same time point.
% \end{remark}

% \subsection{Consequences used later}

% \begin{corollary}[Shape of the rate under linear preparation]
% \label{cor:ramp}
% With $\mu(\tau)=m_0+\alpha\tau$ and $k=\alpha/\sigma_Y$,
% \begin{equation}
% \ln\nu(\tau)=\text{const}-\frac{(a_0-k\tau)^{2}}{2},\qquad
% \frac{\dd}{\dd\tau}\ln\nu=k\,a(\tau),\qquad
% \frac{\dd^{2}}{\dd\tau^{2}}\ln\nu=-k^{2}<0 .
% \end{equation}
% \end{corollary}

% \begin{proof}
% Take logarithms in Proposition~\ref{prop:rate}. The bracket is constant in $\tau$, so it
% contributes only to the constant, and $\ln\varphi(a)=-a^2/2-\frac12\ln2\pi$ with
% $a(\tau)=a_0-k\tau$. Differentiate twice.
% \end{proof}

% The constancy of the bracket is worth noting: the trend correction, being multiplicative and
% level-independent, shifts $\ln\nu$ by a constant and therefore leaves both derivatives
% unchanged. Corollary~\ref{cor:ramp} is exact in the presence of the trend.

% \begin{remark}
% The concavity is parameter-free and is not contradicted by the data. Over the analysed
% window the slope of the observed log rate falls from $+0.809$ (1--2 s) to $+0.702$ (2--3 s),
% a change of $-0.107$ with $p=0.20$, which is the direction Corollary~\ref{cor:ramp}
% requires. A convexity appears only if the window is extended to 0--3 s ($+0.106$,
% $p=0.009$), and that window contains the residual accumulation the manuscript excludes.
% \end{remark}

% \subsection{Where the two anchorings separate}
% \label{sec:anchor-theory}

% \begin{proposition}[Why the anchorings differ on R9]
% \label{prop:R9}
% Let $\bar v$ be the coherence-averaged rate of rise of the mean level while evidence is still
% accumulating. Under offset anchoring, no 2-second trial contributes preparation to the 0--2 s
% window: the level rises at $\bar v$ for uncommitted trials and is flat for committed ones. In
% the 0.3-second condition every trial has committed by 0.52 s and every trial contributes
% preparation, so the level rises at $\alpha$ throughout. The plateau therefore makes the
% 2-second condition rise more slowly on average, which is R9. Under commitment anchoring the
% plateau disappears and 2-second trials prepare inside the window, so R9 requires
% $\alpha>\bar v$.
% \end{proposition}

% The second condition is the one the manuscript states at p.~17, lines 330--331.
% Section~\ref{sec:anchor} reports what happens when it is imposed.

% \subsection{The obstruction to R15}
% \label{sec:obstruction}

% \begin{proposition}[Saturation annihilates the coherence effect]
% \label{prop:sat}
% Let $F_{D,c}$ be the probability of reaching $\theta$ before evidence ends and
% $h(u)=\varphi\bigl((g-u)/\sigma_Y\bigr)/\sigma_Y$, which is strictly increasing on
% $u<g$. Since $|m_0|=\min(|\tilde y|,\theta)$ almost surely, Proposition~\ref{prop:rate}
% gives, at any fixed time in the delay,
% \begin{equation}
% \nu(D,c)=\Gamma_0\Bigl[F_{D,c}\,h(\theta+\alpha\tau)
% +(1-F_{D,c})\,\E\bigl[h(|m_0|+\alpha\tau)\bigm| |m_0|<\theta\bigr]\Bigr],
% \end{equation}
% with $\Gamma_0=\alpha\Phi(m)+\sigma_{\dot Y}\varphi(m)$ independent of $D$ and $c$. If
% $F_{D,c}\to1$ for both coherence levels then $\nu\to\Gamma_0 h(\theta+\alpha\tau)$
% independently of coherence, and the coherence effect vanishes identically.
% \end{proposition}

% \begin{proof}
% Condition on whether the trial reached $\theta$. On the event that it did,
% $|m_0|=\theta$ deterministically; on its complement $|m_0|<\theta$. Take expectations of
% $\nu$ over $m_0$, using Proposition~\ref{prop:rate} and the fact that the bracket does not
% depend on the level. The limit is immediate.
% \end{proof}

% \begin{corollary}
% \label{cor:R15}
% R15 requires the coherence effect to be larger after a 2-second stimulus, hence requires a
% substantial fraction of 2-second trials to fail to reach the internal threshold.
% \end{corollary}

%======================================================================
\section{Method}
\label{sec:method}

\subsection{Analysis windows}

A single module defines the four windows of the manuscript (p.~15, lines 284--289) --- the
stimulus period 0 to 2 s after onset, the delay 1 to 3 s after offset, the response-signal
period 0 to 0.2 s, the response period $-0.2$ to 0 s --- and a single binning function is
applied to data and simulation alike. The dependent variable is the mean number of partial
activations per 0.1 s bin per trial. Trials with a premature response, no response, an
anticipation or a technical problem are excluded, as in the manuscript.

This gives 20678 trials, the number reported. As a check on the extraction, all nineteen
reported effects come out with the correct sign, and the range of motor execution time
across conditions is 3.0 ms against 2.5 ms reported.

\subsection{Calibration and stability}

Ten parameters were calibrated by Nelder--Mead against six accuracies, six premotor
execution times, six motor execution times, six premature rates, six stimulus-period totals
and six delay-period totals. Every reported effect is the mean over six seeds of 3000 trials
per cell, and is called unstable when its mean is smaller than twice its standard deviation
across seeds.

\begin{remark}[Two errors corrected relative to earlier versions]
First, the burst in progress when a trial reached $g+r$ was counted as a partial activation,
because removing the trial from the risk set triggered a spurious end-of-burst event. This
inflated the response-locked windows twenty-fold and inverted R18. Second, earlier
evaluations were run at a gate position where the model produced 0.06 partial activations
per trial against 2.03 observed, so the sign tests were performed on noise. Both are fixed.
\end{remark}

%======================================================================
\section{Results}
\label{sec:simres}

\subsection{Offset anchoring}

Calibrated: $A=4.234$, $\theta=0.893$, $g=1.590$, $r=0.675$, $\kappa=10.78$,
$\sigma_m=1.822$, $\alpha=0.134$, $w=6.16$, $\sigma_q=1.213$, $T_r=0.047$. These imply
$\sigma_Y=0.320$ and $\beta=15.2$ per second.

\begin{table}[htbp]
\centering
\caption{Simulated / observed, offset anchoring, six seeds of 3000 trials per cell.}
\begin{tabular}{lrrrrrr}
\toprule
cell & accuracy & PMET (ms) & MET (ms) & premature & stimulus & delay\\
\midrule
0.3 s, coh 0.00 & 0.500 / 0.493 & 393 / 396 & 153 / 163 & 0.49 / 0.72 & 0.213 / 0.331 & 0.611 / 1.061\\
0.3 s, coh 0.12 & 0.754 / 0.714 & 380 / 384 & 155 / 160 & 0.69 / 0.60 & 0.292 / 0.357 & 0.810 / 1.093\\
0.3 s, coh 0.50 & 0.941 / 0.951 & 342 / 356 & 161 / 160 & 1.46 / 0.97 & 0.507 / 0.411 & 1.393 / 1.146\\
2.0 s, coh 0.00 & 0.498 / 0.504 & 328 / 341 & 164 / 161 & 1.58 / 1.54 & 0.204 / 0.196 & 1.676 / 1.591\\
2.0 s, coh 0.12 & 0.931 / 0.789 & 312 / 326 & 165 / 162 & 1.99 / 1.55 & 0.283 / 0.233 & 1.888 / 1.682\\
2.0 s, coh 0.50 & 0.976 / 0.952 & 312 / 319 & 166 / 162 & 2.01 / 2.72 & 0.344 / 0.330 & 1.914 / 1.894\\
\bottomrule
\end{tabular}
\end{table}

\begin{table}[htbp]
\centering
\caption{All nineteen effects, offset anchoring.}
\label{tab:eff}
\begin{tabular}{lrrrl}
\toprule
& observed & simulated & seed SD & verdict\\
\midrule
R1 & $+0.04990$ & $+0.05395$ & 0.00121 & OK\\
R2 & $+0.03060$ & $+0.03340$ & 0.00167 & OK\\
R3 & $+0.01790$ & $+0.03494$ & 0.00201 & OK\\
R6 & $+0.00566$ & $+0.00301$ & 0.00043 & OK\\
R7 & $+0.00536$ & $+0.01082$ & 0.00064 & OK\\
R8 & $+0.00097$ & $+0.00170$ & 0.00003 & OK\\
R9 & $+0.00064$ & $+0.00061$ & 0.00009 & OK\\
R10 & $+0.00041$ & $+0.00172$ & 0.00018 & OK\\
R12 & $+0.03115$ & $+0.04438$ & 0.00037 & OK\\
R13 & $+0.00970$ & $+0.02550$ & 0.00062 & OK\\
R14 & $+0.00551$ & $+0.00390$ & 0.00008 & OK\\
R15 & $+0.21040$ & $-0.02723$ & 0.00248 & \textbf{FAIL}\\
R16 & $+0.00139$ & $+0.00230$ & 0.00012 & OK\\
R17 & $-0.00032$ & $+0.00207$ & 0.00021 & \textbf{FAIL}\\
R18 & $-0.03602$ & $-0.02408$ & 0.00442 & OK\\
R19 & $+0.03987$ & $+0.05180$ & 0.00320 & OK\\
R20 & $-0.01900$ & $-0.01449$ & 0.00748 & unstable\\
R22 & $-0.02459$ & $-0.02543$ & 0.00172 & OK\\
R23 & $-0.00556$ & $-0.00454$ & 0.00155 & OK\\
\bottomrule
\end{tabular}
\end{table}

R9 is reproduced almost exactly, $+0.00061$ against $+0.00064$. R20 falls below the
stability threshold, which is expected given that the response-signal window has two bins.

\begin{remark}[A systematic bias in magnitude]
The signs are almost all right, but seven effects exceed the observed 95 per cent confidence
interval and none falls below it: R6, R7, R8, R10, R12, R13, R16. The model is uniformly too
reactive to condition. The same direction of error appears in the trial-to-trial variability
of premotor execution time, which the model underestimates by 45 per cent (92--119 ms
simulated against 161--214 ms observed, computed per participant so that between-participant
differences are already removed). A model that understates trial-to-trial variability
mechanically overstates condition effects. This is a single candidate explanation for two
distinct anomalies; it has not been tested.
\end{remark}

\subsection{Commitment anchoring, at its own best fit}
\label{sec:anchor}

Eleven parameters were calibrated for this variant, adding a saturation level $C$ with
$\theta<C<g$, because Proposition~\ref{prop:R9} requires $\alpha>\bar v$ and a large
$\alpha$ without saturation drives the level past the gate during the delay. Best fit:
$A=3.871$, $\theta=0.889$, $g=1.696$, $r=0.683$, $\kappa=13.64$, $\sigma_m=1.765$,
$\alpha=1.198$, $w=6.36$, $\sigma_q=1.214$, $T_r=0.053$, $C=1.189$. Final loss 39.9 against
6.6 for the offset variant.

The calibration does what Proposition~\ref{prop:R9} predicts: it raises $\alpha$ from 0.134
to 1.198, a nine-fold increase, until preparation outruns accumulation. R9 and R6 both
recover the correct sign. But the fit collapses elsewhere.

\begin{table}[htbp]
\centering
\caption{Commitment anchoring at its own best fit: effects that were robust under offset
anchoring and become indeterminate.}
\begin{tabular}{lrrr}
\toprule
& observed & simulated & seed SD\\
\midrule
R1, PMET duration & $+0.04990$ & $+0.00100$ & 0.00115\\
R2, PMET coherence & $+0.03060$ & $+0.00036$ & 0.00180\\
R12, delay duration & $+0.03115$ & $+0.00003$ & 0.00078\\
R13, delay coherence & $+0.00970$ & $-0.00001$ & 0.00090\\
R14, delay ramp & $+0.00551$ & $+0.00001$ & 0.00009\\
R16, delay duration by time & $+0.00139$ & $-0.00003$ & 0.00011\\
\bottomrule
\end{tabular}
\end{table}

Seven of nineteen effects are correct and robust, twelve are indeterminate, and none fails
outright --- because almost every effect has been reduced to zero.

\paragraph{Where the problem comes from.} A chain of three necessities, each forced by the
previous one.

\emph{First}, R9 forces $\alpha>\bar v$ by Proposition~\ref{prop:R9}. With $\alpha=1.198$
the ramp adds $\alpha\times3.22=3.86$ over the delay, far more than $g-\theta=0.81$, so
without a bound the level would cross the gate long before the response signal and every
trial would respond prematurely. A saturation level is therefore forced.

\emph{Second}, $C$ must lie close to $\theta$ to hold premature responses down; the
calibration puts $C-\theta=0.30$.

\emph{Third}, at $\alpha=1.198$ the ramp covers that distance in $0.30/1.198=0.25$ seconds.
Every trial therefore saturates within a quarter second of committing, and saturates at the
\emph{same} value $C$ regardless of coherence and duration. The maintained state is uniform
across all six cells for almost the whole delay, and every effect that depends on
differences in the maintained level is annihilated. The delay-period total falls to 0.884
against 1.411 observed and premature responses to 0.03 per cent against 1.35 per cent.

\paragraph{There is no intermediate window.} Lowering $\alpha$ restores the level differences
but loses R9 and R6; raising it uniformises the state. The two requirements are monotone in
$\alpha$ and opposite in direction.

\begin{remark}
This does not establish that commitment anchoring is wrong in general. The general argument
for it is sound. What it establishes is that \emph{in this design}, where the delay is fixed
at 3 s and the stimulus lasts at most 2 s, the data require the offset. That is a falsifiable
claim, and Section~\ref{sec:thought} gives the experiment that would test it.
\end{remark}

%======================================================================
\section{R15}
\label{sec:R15}

\subsection{The contradiction, stated exactly}

By Corollary~\ref{cor:R15}, R15 requires a substantial fraction of 2-second trials to fail
to reach the internal threshold; otherwise the maintained level saturates at $\theta$ for
every coherence and the coherence effect vanishes. The accuracy data forbid this. Sweeping
$\theta$ with all other parameters at their calibrated values:

\begin{table}[htbp]
\centering
\caption{R15 against accuracy, as a function of the internal threshold.}
\begin{tabular}{lrrr}
\toprule
$\theta$ & R15 & accuracy, 2 s, coh 0.12 & accuracy, 2 s, coh 0.50\\
\midrule
0.89 (calibrated) & $-0.023$ & 0.931 & 0.974\\
1.10 & $-0.025$ & 0.963 & 0.993\\
1.30 & $+0.023$ & 0.975 & 0.997\\
1.50 & $+0.092$ & 0.972 & 0.998\\
1.70 & $+0.103$ & 0.871 & 1.000\\
\midrule
\emph{observed} & $+0.210$ & \emph{0.789} & \emph{0.952}\\
\bottomrule
\end{tabular}
\end{table}

R15 turns positive only for $\theta$ between roughly 1.3 and 1.8, and never reaches half the
observed value. At exactly those values, accuracy at 2 s and 12 per cent coherence rises to
0.97--0.98 against 0.789 observed, and at 50 per cent to 0.997--1.000 against 0.952.

These are the cells from which the manuscript draws its argument that accumulation
terminates. The contradiction is therefore between R15 and that argument, not between R15
and an incidental feature of the fit.

\subsection{Independent confirmation that accumulation terminates}
\label{sec:terminates}

AGPS (2025) Experiment~2 uses three stimulus durations rather than two, which localises the
termination:

\begin{table}[htbp]
\centering
\caption{Accuracy across three stimulus durations, AGPS 2025 Experiment~2.}
\label{tab:acc3}
\begin{tabular}{lrrrrr}
\toprule
coherence & 300 ms & 1000 ms & 2000 ms & gain 300--1000 & gain 1000--2000\\
\midrule
2\% & 0.515 & 0.530 & 0.524 & $+0.015$ ($p=.38$) & $-0.006$ ($p=.69$)\\
11\% & 0.724 & 0.774 & 0.777 & $+0.050$ ($p=.002$) & $+0.003$ ($p=.81$)\\
40\% & 0.965 & 0.964 & 0.969 & $-0.001$ ($p=.90$) & $+0.005$ ($p=.32$)\\
\bottomrule
\end{tabular}
\end{table}

Unbounded accumulation predicts $d'$ ratios of $\sqrt{1000/300}=1.83$ and $\sqrt2=1.41$;
observed at mid coherence are 1.26 then 1.01. Accumulation uses the evidence up to roughly
one second and then stops.

\subsection{Nine mechanisms, none of which works}

Each was implemented, calibrated where relevant, and evaluated on all nineteen effects.

\begin{longtable}{p{0.29\textwidth}p{0.64\textwidth}}
\toprule
mechanism & outcome\\
\midrule
\endhead
Accelerating preparation & Corrects the shape of the rate profile at constant $\alpha$, but
matching total preparation requires doubling $\alpha$, which drives premature responses to
48--58 per cent and MET to 300--320 ms.\\
Force-integration response stage & Reduces premature responses eighteen-fold at matched MET,
a real gain, but does not move R15. The physiologically faithful second-order twitch
performs worse than a first-order leak, and worse still with a distributed motor-unit
pool.\\
Strategic saturation near the gate & Compresses the most strongly prepared trials, which are
the 2-second ones, reducing the coherence effect where R15 requires it increased.\\
Evidence-dependent maintenance noise & Calibration retains $\eta=0.55$, so the direction is
right, but it delivers 40 per cent of the required gradient while raising the loss from 6.6
to 19.6 and collapsing 2-second premature responses to 0.07--0.22 per cent.\\
Saturation level $C$ & Invisible at a 3-second delay for $C>1.33$, so it cannot affect R15
there. It is required for longer delays, for reasons independent of R15.\\
Leak inside the filter & Contracts differences with time constant $1/(\lambda+\mu)<50$ ms,
so it rescales once rather than eroding over seconds; and it does not bound a linearly
growing input, so it does not replace $C$.\\
Leak on the maintained state & Decays toward zero rather than a non-zero endpoint. At
$\mu=0.25$ per second the delay-period rate collapses fourteen-fold.\\
Funnel toward $\pm\theta$ & Uniformises the maintained state: R15 tends to zero without
changing sign and R3 is extinguished with it.\\
Expansive peripheral projection & The only mechanism reproducing R15 at the right magnitude
($+0.223$ against $+0.218$ at $\gamma=2.5$), at the cost of 81 per cent premature responses
and an inverted R3. R15 and the premature rate are monotonically linked through the same
parameter, so no window exists.\\
\bottomrule
\end{longtable}

The pattern is the one Proposition~\ref{prop:sat} predicts: every mechanism that moves R15
acts through the distance $g-m$, and therefore also moves the premature rate, the absolute
rate, or premotor execution time.

\subsection{How robust is R15 itself}

The manuscript reports $F(2,42)=3.474$, $p=.040$, $\eta_p^2=0.142$. A planned contrast on
the extreme coherence levels, computed per participant, gives $+0.2104$, $t(21)=2.08$,
$p=0.0501$, $d=0.44$; Wilcoxon gives $p=0.0738$. Fourteen of twenty-two participants are
positive, sign test $p=0.29$. Under Benjamini--Hochberg correction across the nineteen
effects, R15 does not survive. The between-participant standard deviation is 0.475, that is
2.3 times the mean; the bootstrap interval is $[+0.013,+0.401]$; leave-one-out moves $p$
between 0.016 and 0.094.

At the subject level, the within-subject standard error of R15, obtained by bootstrapping
trials within each participant, has median 0.425 against a between-subject standard
deviation of 0.475. The ratio is 1.12: almost all of the apparent between-subject spread is
sampling noise within participants, so there is very little genuine individual variability
to explain. Compared against the model prediction of $-0.026$, three participants of
twenty-two differ by more than two standard errors, all positive; roughly one would be
expected by chance.

\subsection{R15 in the companion experiment}

AGPS (2025) Experiment~2 reports $F(1.86,39.04)=2.25$, $p=.12$,
$\mathrm{BF}_{\text{incl}}=0.70$, anecdotal evidence for the null.

In matched units --- partial activations per second, extreme coherence contrast --- the two
experiments give $+0.105$ per second ($p=0.0501$, $d=0.44$) in J~Neuroscience and $+0.138$
per second ($p=0.086$, $d=0.38$) in AGPS. The effect is \emph{larger} in the companion
experiment. It replicates in direction and magnitude and fails to reach significance in
either study alone. Stouffer combination gives $z=2.74$, $p=0.006$.

R15 should therefore not be treated as a false positive. It is a modest effect that neither
experiment can establish alone and that the two together do establish. The model's failure
is correspondingly substantive.

\subsection{Critical assessment of the manuscript's treatment}

Three explanations are offered, in two places, and they are mutually inconsistent.

\emph{Exclusion of anticipatory trials} (p.~18, lines 350--359). The manuscript concedes it
is insufficient. It can now be quantified: including premature trials with person-time
denominators changes every reconstructed quantity in the third decimal place, since they are
1.35 per cent of trials and mostly occur late in the delay.

\emph{Strategic reduction of preparation near the gate} (p.~18, lines 359--363). Implemented
and tested, it acts in the wrong direction: it compresses differences among the most
strongly prepared trials, which are the 2-second trials, reducing the coherence effect
exactly where R15 needs it increased. R15 remains negative across the range of preparation
slopes tested.

\emph{Differential threshold reaching} (p.~22--23, lines 460--468). This is the second of
the two branches set out at p.~9--10, lines 143--151. The Results invoke the \emph{first}
branch for R3 --- ``coherence-related differences in the starting point of the second
accumulation phase were smaller in the 2 s condition'' (p.~16, lines 313--315) --- and the
Discussion invokes the second for R15 --- ``increasing motion duration appears to have
amplified coherence-related differences in sustained motor preparatory activity'' (p.~23,
lines 466--468). Both statements concern the same quantity, and the Methods state explicitly
that the branches predict opposite interactions. Simulation confirms they are exclusive: the
calibrated model implements the first branch, which R3 requires, and that forces the wrong
sign for R15.

R24, the apparent levelling off before the response signal, is used to support the
strategic-reduction account. It requires no mechanism: the rate is bounded by $\beta/2\pi$
and its sensitivity $\dd\nu/\dd\mu=\nu a/\sigma_Y$ tends to zero as the level approaches the
gate, so the highest-rate condition flattens first.

\paragraph{Recommendation.} State R15 as a limitation with the contradiction of
Section~\ref{sec:obstruction} made explicit, and withdraw all three proposed explanations.
The contradiction between Results and Discussion should be resolved in one direction or the
other regardless of what is done about R15.

%======================================================================
\section{R3}
\label{sec:R3}

\subsection{Status in the two experiments}

J~Neuroscience: $F(2,42)=3.62$, $p=.035$. AGPS (2025) Experiment~2, nearly the same design:
$F(4,84)<1$, $\mathrm{BF}_{\text{incl}}=0.05$, that is \emph{strong} evidence for the null.

The discrepancy is not a matter of the coherence range at the low end. In J~Neuroscience the
interaction is carried entirely by the 12 versus 50 per cent contrast: dropping the
zero-coherence condition \emph{strengthens} it, from $+0.0179$ ($p=0.075$) to $+0.0210$
($p=0.0044$), while the 0 versus 12 per cent contrast gives $-0.0032$ ($p=0.73$).

\subsection{Where the difference lies}

Simple coherence effects on premotor execution time, mid versus high coherence:

\begin{center}
\begin{tabular}{lrr}
\toprule
& short stimulus & long stimulus\\
\midrule
J~Neuroscience (12 vs 50\%) & $+28.2$ ms ($p=.007$) & $+7.1$ ms ($p=.15$)\\
AGPS Exp.~2 (11 vs 40\%) & $+30.5$ ms & $+27.2$ ms\\
\bottomrule
\end{tabular}
\end{center}

The short-duration effect replicates almost exactly, 28.2 against 30.5 ms. The entire
discrepancy is at long duration: coherence ceases to influence premotor execution time after
2 s in J~Neuroscience and continues to do so in AGPS.

\subsection{The explanation I proposed, and why it is refuted}

The natural explanation is that accumulation terminates before 2 s in one experiment and not
the other, so the maintained level saturates at $\theta$ across coherences in one and not
the other. This is testable through accuracy, since termination implies no further accuracy
gain.

It is refuted. The accuracy gains from short to long stimulus are essentially identical: at
mid coherence $+0.075$ ($p=.0009$) in J~Neuroscience against $+0.053$ ($p=.0146$) in AGPS;
at high coherence $+0.000$ against $+0.004$. The degree of saturation of the decision process
is the same in the two experiments.

\subsection{What this leaves}

At equal decision quality, coherence still influences premotor execution time after a
2-second stimulus in one experiment and not in the other. The difference must therefore lie
between the decision and premotor execution time, not in the accumulation. I have no
explanation.

Two factual differences are on record and neither is obviously sufficient. Stimulus
parameters differ: dot diameter 0.05$^\circ$ against 0.06$^\circ$, dot speed 10.5 against
6.47 degrees per second, and different aperture sizes. And the coherence levels differ at
both ends, 0 and 50 per cent against 2 and 40 per cent.

\subsection{Critical assessment of the manuscript's treatment}

The Methods, p.~9--10, lines 136--156, are careful and correct: the temporal preparation
hypothesis does not predict the interaction; the threshold and combined hypotheses allow it;
and its \emph{direction} is explicitly not predicted, since ``the direction of this
interaction depends on how the additional motion evidence provided by the 2 s duration
changes the relative probability of reaching the internal threshold across coherence
levels''.

That clause has a consequence the manuscript does not draw. If both directions are permitted
in advance, observing one of them discriminates against the temporal preparation hypothesis
but provides no test of the combined hypothesis. R3 should be presented as evidence against
H2, not as evidence for H3.

The account given in the Results also contradicts the account given in the Discussion for
R15, as documented in Section~\ref{sec:R15}.

\paragraph{On robustness.} Given the failure to replicate in the companion experiment, and
given that R3 does not survive multiple-comparison correction, it should be reported with
that context. The model reproduces it ($+0.0349$ simulated against $+0.0179$ observed,
inside the observed confidence interval), but reproducing an effect that does not replicate
is weak support.

%======================================================================
\section{A thought experiment that discriminates the anchorings}
\label{sec:thought}

Consider two conditions with the same total time from stimulus onset to response signal:
condition~1 with a 5-second stimulus and a 5-second delay, condition~2 with an 8-second
stimulus and a 2-second delay.

Both stimuli are long enough that essentially every trajectory reaches $\theta$, so the
maintained level and the accuracy are the same. The only quantity that differs is the
duration of motor preparation.

\begin{proposition}
\label{prop:thought}
Under offset anchoring, preparation runs for the length of the delay, 5 s against 2 s, so
the level at the response signal is higher in condition~1 by $3\alpha$: premotor execution
time is shorter and the delay-period rate higher in condition~1. Under commitment anchoring,
preparation runs from $\Tstop$ to the response signal, $10-\Tstop$ seconds in both, so the
two conditions are predicted to be identical.
\end{proposition}

This requires neither electroencephalography nor measured excursion durations. It is not
realisable in the existing data: no pair of conditions in the three experiments has matched
total time with a different split.

\paragraph{The partial equivalent in the present data.} R1 decomposes differently under the
two anchorings. Under offset anchoring the ramp lasts 3 s in both duration conditions, so the
entire 49.9 ms effect must come from differences in $m_0$. Under commitment anchoring the
ramp lasts 1.71 s longer after a 2-second stimulus, contributing $\alpha\times1.71/w\approx
37$ ms at the offset-calibrated $\alpha$, so most of the effect would come from the ramp.
Both variants can produce a 50 ms effect, so R1 does not discriminate by itself. What
discriminates in the present data is R9 and R6, through the plateau.

\paragraph{An accessible version.} Crossing a short delay with a long stimulus --- 2 s of
stimulus with a 1 s delay against 0.3 s of stimulus with a 1 s delay --- separates the two
without requiring ten-second trials. Under offset anchoring the ramp is 1 s in both and R9
survives; under commitment anchoring the long-stimulus condition prepares throughout the
stimulus and R9 should reverse.

%======================================================================
\section{Is information after the internal threshold discarded?}

The model assumes evidence arriving after $\Tstop$ is discarded entirely. The three stimulus
durations of AGPS Experiment~2 (Table~\ref{tab:acc3}) support the assumption and make it
more precise than the manuscript can: accumulation uses the evidence up to roughly one
second at intermediate coherence and earlier at high coherence, after which further evidence
does not contribute to the decision. The second half of a 2-second stimulus contributes
nothing.

One distinction should be drawn that the manuscript does not. What the accuracy data
establish is that the \emph{choice} is settled, since accuracy depends only on the sign.
They place no constraint on whether the \emph{level} continues to be updated after
commitment. The manuscript treats these as one claim at p.~22, lines 439--443. Separating
them is what the title of the paper implies: a decision that is locked while the motor level
remains graded.

%======================================================================
\section{Summary}

\begin{itemize}
\item With preparation added outside the filter and anchored to stimulus offset, sixteen of
nineteen effects are reproduced robustly, including R9 to within four per cent.
\item Commitment anchoring, the more natural specification and the one originally proposed,
cannot be made to work in this design. R9 forces a nine-fold increase in the preparation
slope; that forces a saturation level; and the saturation uniformises the maintained state,
annihilating six effects the offset variant reproduces.
\item R15 fails, and Proposition~\ref{prop:sat} shows why: it requires 2-second trials not
to reach the internal threshold, which the accuracy data forbid. Nine mechanisms were tested
and none removes the contradiction.
\item R15 is nonetheless a real effect: it replicates in the companion experiment with a
larger effect size, and the two combined give $p=0.006$.
\item R3 is reproduced here but does not replicate in the companion experiment. The
difference lies entirely at long stimulus duration and is not explained by differential
termination of accumulation, since the accuracy gains are the same in the two experiments.
\item The model is uniformly too reactive to condition: seven effects exceed the observed
confidence interval and none falls below. This may be the same deficiency as the 45 per cent
underestimation of trial-to-trial variability in premotor execution time.
\end{itemize}

\end{document}
